\chapter{Reflexive Verbs}

% sich ändern (to change)
\begin{verbentry}{
  style = {acc},
  toc = {sich ändern (to change)},
  name = {sich ändern},
  english = {to change},
  type = {regular, + Akkusativ},
  auxiliary = {haben}
}
 
    
       \vspace{5pt}

    \hrule
    
    \vspace{5pt}

  \begin{tabular}{rl}
     \multicolumn{2}{l}{\textbf{PRÄSENS} }\\
    ich & \textbf{ändere mich}\\
    du & \textbf{änderst dich}\\
    er/sie/es & \textbf{ändert sich}\\
    wir & \textbf{ändern uns}\\
    ihr & \textbf{ändert euch}\\
    sie/Sie & \textbf{ändern sich}\\
  \end{tabular}
  \hfill
     \begin{tabular}{rl}
    \multicolumn{2}{l}{\textbf{PRÄTERITUM} }\\
    ich & \textbf{änderte mich}\\
    du & \textbf{ändertest dich}\\
    er/sie/es & \textbf{änderte sich}\\
    wir & \textbf{änderten uns}\\
    ihr & \textbf{ändertet euch}\\
    sie/Sie & \textbf{änderten sich}\\
  \end{tabular}
  \hfill
  \begin{tabular}{rl}
     \multicolumn{2}{l}{\textbf{IMPERATIV} }\\
   & \textbf{änder(e) (du) dich!}\\
   & \textbf{ändern wir uns!}\\
   & \textbf{ändert (ihr) euch!}\\
   & \textbf{ändern Sie sich!}\\
   & \vspace{10pt}
  \end{tabular}

    \vspace{10pt}

 \begin{tabular}{rl}
     \multicolumn{2}{l}{\textbf{PERFEKT}}\\
    &haben + sich + \textbf{geändert}\\
  \end{tabular}
\hfill
   \begin{tabular}{rl}
   
    \multicolumn{2}{l}{\textbf{FUTURE I} }\\
    &werden + sich + \textbf{ändern}\\
 
  \end{tabular}
\hfill
   \begin{tabular}{rl}
      \multicolumn{2}{l}{\textbf{PLUSQUAMPERFEKT}}\\
    &hatten + sich + \textbf{geändert}\\
  \end{tabular}
  
  
\vspace{10pt}

\begin{tabular}{lll}
\multicolumn{3}{l}{\textbf{MIT PRÄPOSITIONEN}}\\
& \textbf{---} & \\
\end{tabular}
\hfill
\begin{tabular}{lll}
\multicolumn{3}{l}{\textbf{TRENNBARE VERBEN}}\\
& \textbf{---} & \\
\end{tabular}

\vspace{-5pt}
\end{verbentry}

% sich ärgern (to be annoyed)
\begin{verbentry}{
  style = {acc},
  toc = {sich ärgern (to be annoyed)},
  name = {sich ärgern},
  english = {to be annoyed},
  type = {regular, + Akkusativ},
  auxiliary = {haben}
}
 
    
       \vspace{5pt}

    \hrule
    
    \vspace{5pt}

  \begin{tabular}{rl}
     \multicolumn{2}{l}{\textbf{PRÄSENS} }\\
    ich & \textbf{ärgere mich}\\
    du & \textbf{ärgerst dich}\\
    er/sie/es & \textbf{ärgert sich}\\
    wir & \textbf{ärgern uns}\\
    ihr & \textbf{ärgert euch}\\
    sie/Sie & \textbf{ärgern sich}\\
  \end{tabular}
  \hfill
     \begin{tabular}{rl}
    \multicolumn{2}{l}{\textbf{PRÄTERITUM} }\\
    ich & \textbf{ärgerte mich}\\
    du & \textbf{ärgertest dich}\\
    er/sie/es & \textbf{ärgerte sich}\\
    wir & \textbf{ärgerten uns}\\
    ihr & \textbf{ärgertet euch}\\
    sie/Sie & \textbf{ärgerten sich}\\
  \end{tabular}
  \hfill
  \begin{tabular}{rl}
     \multicolumn{2}{l}{\textbf{IMPERATIV} }\\
   & \textbf{ärger(e) (du) dich!}\\
   & \textbf{ärgern wir uns!}\\
   & \textbf{ärgert (ihr) euch!}\\
   & \textbf{ärgern Sie sich!}\\
   & \vspace{10pt}
  \end{tabular}

    \vspace{10pt}

 \begin{tabular}{rl}
     \multicolumn{2}{l}{\textbf{PERFEKT}}\\
    &haben + sich + \textbf{geärgert}\\
  \end{tabular}
\hfill
   \begin{tabular}{rl}
   
    \multicolumn{2}{l}{\textbf{FUTURE I} }\\
    &werden + sich + \textbf{ärgern}\\
 
  \end{tabular}
\hfill
   \begin{tabular}{rl}
      \multicolumn{2}{l}{\textbf{PLUSQUAMPERFEKT}}\\
    &hatten + sich + \textbf{geärgert}\\
  \end{tabular}
  
\vspace{10pt}

   \begin{tabular}{lll}
      \multicolumn{3}{l}{\textbf{MIT PRÄPOSITIONEN}}\\
& \textbf{sich ärgern über} + Akk & (to be annoyed about)
  \end{tabular}
\hfill
\begin{tabular}{lll}
\multicolumn{3}{l}{\textbf{TRENNBARE VERBEN}}\\
& \textbf{---} & \\
\end{tabular}
 % \hfill
 %  \begin{tabular}{lll}
  %    \multicolumn{3}{l}{\textbf{TRENNBARE VERBEN}}\\
%\vspace{10pt}
 % \end{tabular}
  
  \vspace{-5pt}
\end{verbentry}

% sich bedanken (to thank)
\begin{verbentry}{
  style = {default},
  toc = {sich bedanken (to thank)},
  name = {sich bedanken},
  english = {to thank},
  type = {regular},
  auxiliary = {haben}
}
 
    
       \vspace{5pt}

    \hrule
    
    \vspace{5pt}

  \begin{tabular}{rl}
     \multicolumn{2}{l}{\textbf{PRÄSENS} }\\
    ich & \textbf{bedanke mich}\\
    du & \textbf{bedankst dich}\\
    er/sie/es & \textbf{bedankt sich}\\
    wir & \textbf{bedanken uns}\\
    ihr & \textbf{bedankt euch}\\
    sie/Sie & \textbf{bedanken sich}\\
  \end{tabular}
  \hfill
     \begin{tabular}{rl}
    \multicolumn{2}{l}{\textbf{PRÄTERITUM} }\\
    ich & \textbf{bedankte mich}\\
    du & \textbf{bedanktest dich}\\
    er/sie/es & \textbf{vbedankte sich}\\
    wir & \textbf{bedankten uns}\\
    ihr & \textbf{bedanktet euch}\\
    sie/Sie & \textbf{bedankten sich}\\
  \end{tabular}
  \hfill
  \begin{tabular}{rl}
     \multicolumn{2}{l}{\textbf{IMPERATIV} }\\
   & \textbf{bedank(e) (du) dich!}\\
   & \textbf{bedanken wir uns!}\\
   & \textbf{bedankt (ihr) euch!}\\
   & \textbf{bedanken Sie sich!}\\
   & \vspace{10pt}
  \end{tabular}

    \vspace{10pt}

 \begin{tabular}{rl}
     \multicolumn{2}{l}{\textbf{PERFEKT}}\\
    &haben + sich + \textbf{bedankt}\\
  \end{tabular}
\hfill
   \begin{tabular}{rl}
   
    \multicolumn{2}{l}{\textbf{FUTURE I} }\\
    &werden + sich + \textbf{bedanken}\\
 
  \end{tabular}
\hfill
   \begin{tabular}{rl}
      \multicolumn{2}{l}{\textbf{PLUSQUAMPERFEKT}}\\
    &hatten + sich + \textbf{bedankt}\\
  \end{tabular}
  
\vspace{10pt}

   \begin{tabular}{lll}
      \multicolumn{3}{l}{\textbf{MIT PRÄPOSITIONEN}}\\
& \textbf{sich bedanken für} + Akk & (to thank for)\\
& \textbf{sich bedanken bei} + Dat & (to be thankful to someone)
  \end{tabular}
 \hfill
  \begin{tabular}{lll}
\multicolumn{3}{l}{\textbf{TRENNBARE VERBEN}}\\
& \textbf{---} & \\
\end{tabular}
  
  \vspace{-5pt}
\end{verbentry}

% sich beeilen (to hurry)
\begin{verbentry}{
  style = {default},
  toc = {sich beeilen (to hurry)},
  name = {sich beeilen},
  english = {to hurry},
  type = {regular},
  auxiliary = {haben}
}
 
    
       \vspace{5pt}

    \hrule
    
    \vspace{5pt}

  \begin{tabular}{rl}
     \multicolumn{2}{l}{\textbf{PRÄSENS} }\\
    ich & \textbf{beeile mich}\\
    du & \textbf{beeilst dich}\\
    er/sie/es & \textbf{beeilt sich}\\
    wir & \textbf{beeilen uns}\\
    ihr & \textbf{beeilt euch}\\
    sie/Sie & \textbf{beeilen sich}\\
  \end{tabular}
  \hfill
     \begin{tabular}{rl}
    \multicolumn{2}{l}{\textbf{PRÄTERITUM} }\\
    ich & \textbf{beeilte mich}\\
    du & \textbf{beeiltest dich}\\
    er/sie/es & \textbf{beeilte sich}\\
    wir & \textbf{beeilten uns}\\
    ihr & \textbf{beeiltet euch}\\
    sie/Sie & \textbf{beeilten sich}\\
  \end{tabular}
  \hfill
  \begin{tabular}{rl}
     \multicolumn{2}{l}{\textbf{IMPERATIV} }\\
   & \textbf{beeil(e) (du) dich!}\\
   & \textbf{beeilen wir uns!}\\
   & \textbf{beeilt (ihr) euch!}\\
   & \textbf{beeilen Sie sich!}\\
   & \vspace{10pt}
  \end{tabular}

    \vspace{10pt}

 \begin{tabular}{rl}
     \multicolumn{2}{l}{\textbf{PERFEKT}}\\
    &haben + sich + \textbf{beeilt}\\
  \end{tabular}
\hfill
   \begin{tabular}{rl}
   
    \multicolumn{2}{l}{\textbf{FUTURE I} }\\
    &werden + sich + \textbf{beeilen}\\
 
  \end{tabular}
\hfill
   \begin{tabular}{rl}
      \multicolumn{2}{l}{\textbf{PLUSQUAMPERFEKT}}\\
    &hatten + sich + \textbf{beeilt}\\
  \end{tabular}
  
\vspace{10pt}

   \begin{tabular}{lll}
      \multicolumn{3}{l}{\textbf{MIT PRÄPOSITIONEN}}\\
& \textbf{sich beeilen mit} + Dat & (to hurry up with)
  \end{tabular}
\hfill
\begin{tabular}{lll}
\multicolumn{3}{l}{\textbf{TRENNBARE VERBEN}}\\
& \textbf{---} & \\
\end{tabular}
 % \hfill
 %  \begin{tabular}{lll}
  %    \multicolumn{3}{l}{\textbf{TRENNBARE VERBEN}}\\
%\vspace{10pt}
 % \end{tabular}
  
  \vspace{-5pt}
\end{verbentry}

% sich befinden (to be located)
\begin{verbentry}{
  style = {default},
  toc = {sich befinden (to be located)},
  name = {sich befinden},
  english = {to be located},
  type = {irregular},
  auxiliary = {haben}
}
 
    
       \vspace{5pt}

    \hrule
    
    \vspace{5pt}

  \begin{tabular}{rl}
     \multicolumn{2}{l}{\textbf{PRÄSENS} }\\
    ich & \textbf{befinde mich}\\
    du & \textbf{befindest dich}\\
    er/sie/es & \textbf{befindet sich}\\
    wir & \textbf{befinden uns}\\
    ihr & \textbf{befindet euch}\\
    sie/Sie & \textbf{befinden sich}\\
  \end{tabular}
  \hfill
     \begin{tabular}{rl}
    \multicolumn{2}{l}{\textbf{PRÄTERITUM} }\\
    ich & \textbf{befand mich}\\
    du & \textbf{befandest dich}\\
    er/sie/es & \textbf{befand sich}\\
    wir & \textbf{befanden uns}\\
    ihr & \textbf{befandet euch}\\
    sie/Sie & \textbf{befanden sich}\\
  \end{tabular}
  \hfill
  \begin{tabular}{rl}
     \multicolumn{2}{l}{\textbf{IMPERATIV} }\\
   & \textbf{befind(e) (du) dich!}\\
   & \textbf{befinden wir uns!}\\
   & \textbf{befindet (ihr) euch!}\\
   & \textbf{befinden Sie sich!}\\
   & \vspace{10pt}
  \end{tabular}

    \vspace{10pt}

 \begin{tabular}{rl}
     \multicolumn{2}{l}{\textbf{PERFEKT}}\\
    &haben + sich + \textbf{befunden}\\
  \end{tabular}
\hfill
   \begin{tabular}{rl}
   
    \multicolumn{2}{l}{\textbf{FUTURE I} }\\
    &werden + sich + \textbf{befinden}\\
 
  \end{tabular}
\hfill
   \begin{tabular}{rl}
      \multicolumn{2}{l}{\textbf{PLUSQUAMPERFEKT}}\\
    &hatten + sich + \textbf{befunden}\\
  \end{tabular}
  
\vspace{10pt}

   \begin{tabular}{lll}
      \multicolumn{3}{l}{\textbf{MIT PRÄPOSITIONEN}}\\
& \textbf{sich befinden in} + Akk & (to be located in)
  \end{tabular}
  \hfill
   \begin{tabular}{lll}
\multicolumn{3}{l}{\textbf{TRENNBARE VERBEN}}\\
& \textbf{---} & \\
\end{tabular}
  
  \vspace{-5pt}
\end{verbentry}

% sich bemühen (to make an effort)
\begin{verbentry}{
  style = {default},
  toc = {sich bemühen (to make an effort)},
  name = {sich bemühen},
  english = {to make an effort},
  type = {regular, reflexive},
  auxiliary = {haben}
}
 
    
     \vspace{5pt}

    \hrule
    
    \vspace{5pt}

  \begin{tabular}{rl}
     \multicolumn{2}{l}{\textbf{PRÄSENS} }\\
    ich & \textbf{bemühe mich}\\
    du & \textbf{bemühst dich}\\
    er/sie/es & \textbf{bemüht sich}\\
    wir & \textbf{bemühen uns}\\
    ihr & \textbf{bemüht euch}\\
    sie/Sie & \textbf{bemühen sich}\\
  \end{tabular}
  \hfill
     \begin{tabular}{rl}
    \multicolumn{2}{l}{\textbf{PRÄTERITUM} }\\
    ich & \textbf{bemühte mich}\\
    du & \textbf{bemühtest dich}\\
    er/sie/es & \textbf{bemühte sich}\\
    wir & \textbf{bemühten uns}\\
    ihr & \textbf{bemühtet euch}\\
    sie/Sie & \textbf{bemühten sich}\\
  \end{tabular}
  \hfill
  \begin{tabular}{rl}
     \multicolumn{2}{l}{\textbf{IMPERATIV} }\\
   & \textbf{bemüh(e) dich!}\\
   & \textbf{bemühen wir uns!}\\
   & \textbf{bemüht euch!}\\
   & \textbf{bemühen Sie sich!}\\
   & \vspace{10pt}
  \end{tabular}

    \vspace{10pt}

 \begin{tabular}{rl}
     \multicolumn{2}{l}{\textbf{PERFEKT}}\\
   & haben + sich + \textbf{bemüht}\\
  \end{tabular}
\hfill
   \begin{tabular}{rl}
   
    \multicolumn{2}{l}{\textbf{FUTURE I} }\\
   & werden + sich + \textbf{bemühen}\\
 
  \end{tabular}
\hfill
   \begin{tabular}{rl}
      \multicolumn{2}{l}{\textbf{PLUSQUAMPERFEKT}}\\
   & hatten + sich + \textbf{bemüht}\\
  \end{tabular}

   \vspace{10pt}

   \begin{tabular}{lll}
      \multicolumn{3}{l}{\textbf{MIT PRÄPOSITIONEN}}\\
& \textbf{sich bemühen um} + Akk & (to make an effort for / to seek)\\
& \textbf{sich bemühen zu} + Inf & (to try to)\\
\end{tabular}
  \hfill
   \begin{tabular}{lll}
\multicolumn{3}{l}{\textbf{TRENNBARE VERBEN}}\\
& \textbf{---} & \\
\end{tabular}


  \vspace{-5pt}
\end{verbentry}

% sich beschäftigen (to occupy oneself / to deal with)
\begin{verbentry}{
  style = {default},
  toc = {sich beschäftigen (to occupy oneself / to deal with)},
  name = {sich beschäftigen},
  english = {to occupy oneself / to deal with},
  type = {regular},
  auxiliary = {haben}
}
 
    
     \vspace{5pt}

    \hrule
    
    \vspace{5pt}

  \begin{tabular}{rl}
     \multicolumn{2}{l}{\textbf{PRÄSENS} }\\
    ich & \textbf{beschäftige mich}\\
    du & \textbf{beschäftigst dich}\\
    er/sie/es & \textbf{beschäftigt sich}\\
    wir & \textbf{beschäftigen uns}\\
    ihr & \textbf{beschäftigt euch}\\
    sie/Sie & \textbf{beschäftigen sich}\\
  \end{tabular}
  \hfill
     \begin{tabular}{rl}
    \multicolumn{2}{l}{\textbf{PRÄTERITUM} }\\
    ich & \textbf{beschäftigte mich}\\
    du & \textbf{beschäftigtest dich}\\
    er/sie/es & \textbf{beschäftigte sich}\\
    wir & \textbf{beschäftigten uns}\\
    ihr & \textbf{beschäftigtet euch}\\
    sie/Sie & \textbf{beschäftigten sich}\\
  \end{tabular}
  \hfill
  \begin{tabular}{rl}
     \multicolumn{2}{l}{\textbf{IMPERATIV} }\\
   & \textbf{beschäftige dich!}\\
   & \textbf{beschäftigen wir uns!}\\
   & \textbf{beschäftigt euch!}\\
   & \textbf{beschäftigen Sie sich!}\\
   & \vspace{10pt}
  \end{tabular}

    \vspace{10pt}

 \begin{tabular}{rl}
     \multicolumn{2}{l}{\textbf{PERFEKT}}\\
   & haben + sich + \textbf{beschäftigt}\\
  \end{tabular}
\hfill
   \begin{tabular}{rl}
    \multicolumn{2}{l}{\textbf{FUTURE I} }\\
   & werden + sich + \textbf{beschäftigen}\\
  \end{tabular}
\hfill
   \begin{tabular}{rl}
      \multicolumn{2}{l}{\textbf{PLUSQUAMPERFEKT}}\\
   & hatten + sich + \textbf{beschäftigt}\\
  \end{tabular}

   \vspace{10pt}

   \begin{tabular}{lll}
      \multicolumn{3}{l}{\textbf{MIT PRÄPOSITIONEN}}\\
& \textbf{sich beschäftigen mit} + Dat & (to deal with / be occupied with)\\
\end{tabular}
  \hfill
   \begin{tabular}{lll}
\multicolumn{3}{l}{\textbf{TRENNBARE VERBEN}}\\
& \textbf{---} & \\
\end{tabular}

  \vspace{-5pt}
\end{verbentry}

% sich beschweren (to complain)
\begin{verbentry}{
  style = {default},
  toc = {sich beschweren (to complain)},
  name = {sich beschweren},
  english = {to complain},
  type = {regular},
  auxiliary = {haben}
}
 
    
       \vspace{5pt}

    \hrule
    
    \vspace{5pt}

  \begin{tabular}{rl}
     \multicolumn{2}{l}{\textbf{PRÄSENS} }\\
    ich & \textbf{beschwere mich}\\
    du & \textbf{beschwerst dich}\\
    er/sie/es & \textbf{beschwert sich}\\
    wir & \textbf{beschweren uns}\\
    ihr & \textbf{beschwert euch}\\
    sie/Sie & \textbf{beschweren sich}\\
  \end{tabular}
  \hfill
     \begin{tabular}{rl}
    \multicolumn{2}{l}{\textbf{PRÄTERITUM} }\\
    ich & \textbf{beschwerte mich}\\
    du & \textbf{beschwertest dich}\\
    er/sie/es & \textbf{beschwerte sich}\\
    wir & \textbf{beschwerten uns}\\
    ihr & \textbf{beschwertet euch}\\
    sie/Sie & \textbf{beschwerten sich}\\
  \end{tabular}
  \hfill
  \begin{tabular}{rl}
     \multicolumn{2}{l}{\textbf{IMPERATIV} }\\
   & \textbf{beschwer(e) (du) dich!}\\
   & \textbf{beschweren wir uns!}\\
   & \textbf{beschwert (ihr) euch!}\\
   & \textbf{beschweren Sie sich!}\\
   & \vspace{10pt}
  \end{tabular}

    \vspace{10pt}

 \begin{tabular}{rl}
     \multicolumn{2}{l}{\textbf{PERFEKT}}\\
    &haben + sich + \textbf{beschwert}\\
  \end{tabular}
\hfill
   \begin{tabular}{rl}
   
    \multicolumn{2}{l}{\textbf{FUTURE I} }\\
    &werden + sich + \textbf{beschweren}\\
 
  \end{tabular}
\hfill
   \begin{tabular}{rl}
      \multicolumn{2}{l}{\textbf{PLUSQUAMPERFEKT}}\\
    &hatten + sich + \textbf{beschwert}\\
  \end{tabular}
  
\vspace{10pt}

   \begin{tabular}{lll}
      \multicolumn{3}{l}{\textbf{MIT PRÄPOSITIONEN}}\\
& \textbf{sich beschweren über} + Akk & (to complain about something)
  \end{tabular}
  \hfill
   \begin{tabular}{lll}
\multicolumn{3}{l}{\textbf{TRENNBARE VERBEN}}\\
& \textbf{---} & \\
\end{tabular}
  
  \vspace{-5pt}
\end{verbentry}

% sich bewegen (to move oneself)
\begin{verbentry}{
  style = {default},
  toc = {sich bewegen (to move oneself)},
  name = {sich bewegen},
  english = {to move oneself},
  type = {irregular},
  auxiliary = {haben / sein}
}
 
    
     \vspace{5pt}
    \hrule
    \vspace{5pt}

  \begin{tabular}{rl}
     \multicolumn{2}{l}{\textbf{PRÄSENS} }\\
    ich & \textbf{bewege mich}\\
    du & \textbf{bewegst dich}\\
    er/sie/es & \textbf{bewegt sich}\\
    wir & \textbf{bewegen uns}\\
    ihr & \textbf{bewegt euch}\\
    sie/Sie & \textbf{bewegen sich}\\
  \end{tabular}
  \hfill
     \begin{tabular}{rl}
    \multicolumn{2}{l}{\textbf{PRÄTERITUM} }\\
    ich & \textbf{bewegte mich}\\
    du & \textbf{bewegtest dich}\\
    er/sie/es & \textbf{bewegte sich}\\
    wir & \textbf{bewegten uns}\\
    ihr & \textbf{bewegtet euch}\\
    sie/Sie & \textbf{bewegten sich}\\
  \end{tabular}
  \hfill
  \begin{tabular}{rl}
     \multicolumn{2}{l}{\textbf{IMPERATIV} }\\
   & \textbf{beweg dich (du)!}\\
   & \textbf{bewegen wir uns!}\\
   & \textbf{bewegt euch (ihr)!}\\
   & \textbf{bewegen Sie sich!}\\
   & \vspace{10pt}
  \end{tabular}

 \vspace{10pt}

 \begin{tabular}{rl}
     \multicolumn{2}{l}{\textbf{PERFEKT}}\\
   & haben / sein + sich + \textbf{bewegt}
  \end{tabular}
\hfill
   \begin{tabular}{rl}
    \multicolumn{2}{l}{\textbf{FUTURE I}}\\
   & werden + sich + \textbf{bewegen}\\
  \end{tabular}
\hfill
   \begin{tabular}{rl}
      \multicolumn{2}{l}{\textbf{PLUSQUAMPERFEKT}}\\
   & haben / waren + sich + \textbf{bewegt}\\
  \end{tabular}

   \vspace{10pt}

   \begin{tabular}{lll}
      \multicolumn{3}{l}{\textbf{MIT PRÄPOSITIONEN}}\\
& \textbf{sich bewegen in} + Dat & (to move in)\\
& \textbf{sich bewegen auf} + Akk & (to move toward)\\
& \textbf{sich bewegen durch} + Akk & (to move through)\\
\end{tabular}
  \hfill
\begin{tabular}{lll}
\multicolumn{3}{l}{\textbf{TRENNBARE VERBEN}}\\
& \textbf{---} & \\
\end{tabular}

  \vspace{-5pt}
\end{verbentry}

% sich drehen (to turn / to rotate)
\begin{verbentry}{
  style = {default},
  toc = {sich drehen (to turn / to rotate)},
  name = {sich drehen},
  english = {to turn / to rotate},
  type = {regular},
  auxiliary = {haben}
}
 
    
     \vspace{5pt}

    \hrule
    
    \vspace{5pt}

  \begin{tabular}{rl}
     \multicolumn{2}{l}{\textbf{PRÄSENS} }\\
    ich & \textbf{drehe mich}\\
    du & \textbf{drehst dich}\\
    er/sie/es & \textbf{dreht sich}\\
    wir & \textbf{drehen uns}\\
    ihr & \textbf{dreht euch}\\
    sie/Sie & \textbf{drehen sich}\\
  \end{tabular}
  \hfill
     \begin{tabular}{rl}
    \multicolumn{2}{l}{\textbf{PRÄTERITUM} }\\
    ich & \textbf{drehte mich}\\
    du & \textbf{drehtest dich}\\
    er/sie/es & \textbf{drehte sich}\\
    wir & \textbf{drehten uns}\\
    ihr & \textbf{drehtet euch}\\
    sie/Sie & \textbf{drehten sich}\\
  \end{tabular}
  \hfill
  \begin{tabular}{rl}
     \multicolumn{2}{l}{\textbf{IMPERATIV} }\\
   & \textbf{dreh(e) dich!}\\
   & \textbf{drehen wir uns!}\\
   & \textbf{dreht euch!}\\
   & \textbf{drehen Sie sich!}\\
   & \vspace{10pt}
  \end{tabular}

    \vspace{10pt}

 \begin{tabular}{rl}
     \multicolumn{2}{l}{\textbf{PERFEKT}}\\
   & haben + sich + \textbf{gedreht}\\
  \end{tabular}
\hfill
   \begin{tabular}{rl}
    \multicolumn{2}{l}{\textbf{FUTURE I} }\\
   & werden + sich + \textbf{drehen}\\
  \end{tabular}
\hfill
   \begin{tabular}{rl}
      \multicolumn{2}{l}{\textbf{PLUSQUAMPERFEKT}}\\
   & hatten + sich + \textbf{gedreht}\\
  \end{tabular}

   \vspace{10pt}

   \begin{tabular}{lll}
      \multicolumn{3}{l}{\textbf{MIT PRÄPOSITIONEN}}\\
& \textbf{sich drehen um} + Akk & (to revolve around / be about)\\
& \textbf{sich drehen nach} + Dat & (to turn toward)\\
\end{tabular}
  \hfill
   \begin{tabular}{lll}
      \multicolumn{3}{l}{\textbf{TRENNBARE VERBEN}}\\
 & \textbf{um$|$drehen} & (to turn around)\\
 & \textbf{ab$|$drehen} & (to turn off / shoot a film)\\
\vspace{10pt}
  \end{tabular}

  \vspace{-5pt}
\end{verbentry}

% sich eignen (to be suitable)
\begin{verbentry}{
  style = {default},
  toc = {sich eignen (to be suitable)},
  name = {sich eignen},
  english = {to be suitable},
  type = {regular},
  auxiliary = {haben}
}
 
    
     \vspace{5pt}

    \hrule
    
    \vspace{5pt}

  \begin{tabular}{rl}
     \multicolumn{2}{l}{\textbf{PRÄSENS} }\\
    ich & \textbf{eigne mich}\\
    du & \textbf{eignest dich}\\
    er/sie/es & \textbf{eignet sich}\\
    wir & \textbf{eignen uns}\\
    ihr & \textbf{eignet euch}\\
    sie/Sie & \textbf{eignen sich}\\
  \end{tabular}
  \hfill
     \begin{tabular}{rl}
    \multicolumn{2}{l}{\textbf{PRÄTERITUM} }\\
    ich & \textbf{eignete mich}\\
    du & \textbf{eignetest dich}\\
    er/sie/es & \textbf{eignete sich}\\
    wir & \textbf{eigneten uns}\\
    ihr & \textbf{eignetet euch}\\
    sie/Sie & \textbf{eigneten sich}\\
  \end{tabular}
  \hfill
  \begin{tabular}{rl}
     \multicolumn{2}{l}{\textbf{IMPERATIV} }\\
   & \textbf{eigne dich!}\\
   & \textbf{eignen wir uns!}\\
   & \textbf{eignet euch!}\\
   & \textbf{eignen Sie sich!}\\
   & \vspace{10pt}
  \end{tabular}

    \vspace{10pt}

 \begin{tabular}{rl}
     \multicolumn{2}{l}{\textbf{PERFEKT}}\\
   & haben + sich + \textbf{geeignet}\\
  \end{tabular}
\hfill
   \begin{tabular}{rl}
    \multicolumn{2}{l}{\textbf{FUTURE I} }\\
   & werden + sich + \textbf{eignen}\\
  \end{tabular}
\hfill
   \begin{tabular}{rl}
      \multicolumn{2}{l}{\textbf{PLUSQUAMPERFEKT}}\\
   & hatten + sich + \textbf{geeignet}\\
  \end{tabular}

   \vspace{10pt}

   \begin{tabular}{lll}
      \multicolumn{3}{l}{\textbf{MIT PRÄPOSITIONEN}}\\
& \textbf{sich eignen für} + Akk & (to be suitable for)\\
& \textbf{sich eignen als} + Nom & (to be suitable as)\\
\end{tabular}
  \hfill
   \begin{tabular}{lll}
\multicolumn{3}{l}{\textbf{TRENNBARE VERBEN}}\\
& \textbf{---} & \\
\end{tabular}

  \vspace{-5pt}
\end{verbentry}

% sich erinnern (to remember)
\begin{verbentry}{
  style = {default},
  toc = {sich erinnern (to remember)},
  name = {sich erinnern},
  english = {to remember},
  type = {regular},
  auxiliary = {haben}
}
 
    
       \vspace{5pt}

    \hrule
    
    \vspace{5pt}

  \begin{tabular}{rl}
     \multicolumn{2}{l}{\textbf{PRÄSENS} }\\
    ich & \textbf{erinnere mich}\\
    du & \textbf{erinnerst dich}\\
    er/sie/es & \textbf{erinnert sich}\\
    wir & \textbf{erinnern uns}\\
    ihr & \textbf{erinnert euch}\\
    sie/Sie & \textbf{erinnern sich}\\
  \end{tabular}
  \hfill
     \begin{tabular}{rl}
    \multicolumn{2}{l}{\textbf{PRÄTERITUM} }\\
    ich & \textbf{erinnerte mich}\\
    du & \textbf{erinnertest dich}\\
    er/sie/es & \textbf{erinnerte sich}\\
    wir & \textbf{erinnerten uns}\\
    ihr & \textbf{erinnertet euch}\\
    sie/Sie & \textbf{erinnerten sich}\\
  \end{tabular}
  \hfill
  \begin{tabular}{rl}
     \multicolumn{2}{l}{\textbf{IMPERATIV} }\\
   & \textbf{erinner(e) (du) dich!}\\
   & \textbf{erinnern wir uns!}\\
   & \textbf{erinnert (ihr) euch!}\\
   & \textbf{erinnern Sie sich!}\\
   & \vspace{10pt}
  \end{tabular}

    \vspace{10pt}

 \begin{tabular}{rl}
     \multicolumn{2}{l}{\textbf{PERFEKT}}\\
    &haben + sich + \textbf{erinnert}\\
  \end{tabular}
\hfill
   \begin{tabular}{rl}
   
    \multicolumn{2}{l}{\textbf{FUTURE I} }\\
    &werden + sich + \textbf{erinnern}\\
 
  \end{tabular}
\hfill
   \begin{tabular}{rl}
      \multicolumn{2}{l}{\textbf{PLUSQUAMPERFEKT}}\\
    &hatten + sich + \textbf{erinnert}\\
  \end{tabular}
  
\vspace{10pt}

   \begin{tabular}{lll}
      \multicolumn{3}{l}{\textbf{MIT PRÄPOSITIONEN}}\\
& \textbf{sich erinnern an} + Akk & (to recall something)
  \end{tabular}
  \hfill
   \begin{tabular}{lll}
\multicolumn{3}{l}{\textbf{TRENNBARE VERBEN}}\\
& \textbf{---} & \\
\end{tabular}
  
  \vspace{-5pt}
\end{verbentry}

% sich erkundigen (to inquire)
\begin{verbentry}{
  style = {default},
  toc = {sich erkundigen (to inquire)},
  name = {sich erkundigen},
  english = {to inquire},
  type = {regular},
  auxiliary = {haben}
}
 
    
     \vspace{5pt}

    \hrule
    
    \vspace{5pt}

  \begin{tabular}{rl}
     \multicolumn{2}{l}{\textbf{PRÄSENS} }\\
    ich & \textbf{erkundige mich}\\
    du & \textbf{erkundigst dich}\\
    er/sie/es & \textbf{erkundigt sich}\\
    wir & \textbf{erkundigen uns}\\
    ihr & \textbf{erkundigt euch}\\
    sie/Sie & \textbf{erkundigen sich}\\
  \end{tabular}
  \hfill
     \begin{tabular}{rl}
    \multicolumn{2}{l}{\textbf{PRÄTERITUM} }\\
    ich & \textbf{erkundigte mich}\\
    du & \textbf{erkundigtest dich}\\
    er/sie/es & \textbf{erkundigte sich}\\
    wir & \textbf{erkundigten uns}\\
    ihr & \textbf{erkundigtet euch}\\
    sie/Sie & \textbf{erkundigten sich}\\
  \end{tabular}
  \hfill
  \begin{tabular}{rl}
     \multicolumn{2}{l}{\textbf{IMPERATIV} }\\
   & \textbf{erkundige dich!}\\
   & \textbf{erkundigen wir uns!}\\
   & \textbf{erkundigt euch!}\\
   & \textbf{erkundigen Sie sich!}\\
   & \vspace{10pt}
  \end{tabular}

    \vspace{10pt}

 \begin{tabular}{rl}
     \multicolumn{2}{l}{\textbf{PERFEKT}}\\
   & haben + sich + \textbf{erkundigt}\\
  \end{tabular}
\hfill
   \begin{tabular}{rl}
    \multicolumn{2}{l}{\textbf{FUTURE I} }\\
   & werden + sich + \textbf{erkundigen}\\
  \end{tabular}
\hfill
   \begin{tabular}{rl}
      \multicolumn{2}{l}{\textbf{PLUSQUAMPERFEKT}}\\
   & hatten + sich + \textbf{erkundigt}\\
  \end{tabular}

   \vspace{10pt}

   \begin{tabular}{lll}
      \multicolumn{3}{l}{\textbf{MIT PRÄPOSITIONEN}}\\
& \textbf{sich erkundigen nach} + Dat & (to inquire about)\\
& \textbf{sich erkundigen bei} + Dat & (to inquire with / ask at)\\
\end{tabular}
  \hfill
   \begin{tabular}{lll}
\multicolumn{3}{l}{\textbf{TRENNBARE VERBEN}}\\
& \textbf{---} & \\
\end{tabular}

  \vspace{-5pt}
\end{verbentry}

% sich ernähren (to eat / to nourish oneself)
\begin{verbentry}{
  style = {default},
  toc = {sich ernähren (to eat / to nourish oneself)},
  name = {sich ernähren},
  english = {to eat / nourish oneself},
  type = {regular},
  auxiliary = {haben}
}
 
    
     \vspace{5pt}
    \hrule
    \vspace{5pt}

  \begin{tabular}{rl}
     \multicolumn{2}{l}{\textbf{PRÄSENS} }\\
    ich & \textbf{ernähre mich}\\
    du & \textbf{ernährst dich}\\
    er/sie/es & \textbf{ernährt sich}\\
    wir & \textbf{ernähren uns}\\
    ihr & \textbf{ernährt euch}\\
    sie/Sie & \textbf{ernähren sich}\\
  \end{tabular}
  \hfill
     \begin{tabular}{rl}
    \multicolumn{2}{l}{\textbf{PRÄTERITUM} }\\
    ich & \textbf{ernährte mich}\\
    du & \textbf{ernährtest dich}\\
    er/sie/es & \textbf{ernährte sich}\\
    wir & \textbf{ernährten uns}\\
    ihr & \textbf{ernährtet euch}\\
    sie/Sie & \textbf{ernährten sich}\\
  \end{tabular}
  \hfill
  \begin{tabular}{rl}
     \multicolumn{2}{l}{\textbf{IMPERATIV} }\\
   & \textbf{ernähr dich (du)!}\\
   & \textbf{ernähren wir uns!}\\
   & \textbf{ernährt euch (ihr)!}\\
   & \textbf{ernähren Sie sich!}\\
   & \vspace{10pt}
  \end{tabular}

 \vspace{10pt}

 \begin{tabular}{rl}
     \multicolumn{2}{l}{\textbf{PERFEKT}}\\
   & haben + sich + \textbf{ernährt}\\
  \end{tabular}
\hfill
   \begin{tabular}{rl}
    \multicolumn{2}{l}{\textbf{FUTURE I}}\\
   & werden + sich + \textbf{ernähren}\\
  \end{tabular}
\hfill
   \begin{tabular}{rl}
      \multicolumn{2}{l}{\textbf{PLUSQUAMPERFEKT}}\\
   & hatten + sich + \textbf{ernährt}\\
  \end{tabular}

   \vspace{10pt}

   \begin{tabular}{lll}
      \multicolumn{3}{l}{\textbf{MIT PRÄPOSITIONEN}}\\
& \textbf{sich ernähren von} + Dat & (to live on / eat mainly)\\
& \textbf{sich ernähren mit} + Dat & (to nourish oneself with)\\
\end{tabular}
  \hfill
\begin{tabular}{lll}
\multicolumn{3}{l}{\textbf{TRENNBARE VERBEN}}\\
& \textbf{---} & \\
\end{tabular}
\end{verbentry}

% sich engagieren (to get involved / commit oneself)
\begin{verbentry}{
  style = {default},
  toc = {sich engagieren (to get involved / commit oneself)},
  name = {sich engagieren},
  english = {to get involved / commit oneself},
  type = {regular},
  auxiliary = {haben}
}


\vspace{5pt}
\hrule
\vspace{5pt}

\begin{tabular}{rl}
\multicolumn{2}{l}{\textbf{PRÄSENS}}\\
ich & \textbf{engagiere mich}\\
du & \textbf{engagierst dich}\\
er/sie/es & \textbf{engagiert sich}\\
wir & \textbf{engagieren uns}\\
ihr & \textbf{engagiert euch}\\
sie/Sie & \textbf{engagieren sich}\\
\end{tabular}
\hfill
\begin{tabular}{rl}
\multicolumn{2}{l}{\textbf{PRÄTERITUM}}\\
ich & \textbf{engagierte mich}\\
du & \textbf{engagiertest dich}\\
er/sie/es & \textbf{engagierte sich}\\
wir & \textbf{engagierten uns}\\
ihr & \textbf{engagiertet euch}\\
sie/Sie & \textbf{engagierten sich}\\
\end{tabular}
\hfill
\begin{tabular}{rl}
\multicolumn{2}{l}{\textbf{IMPERATIV}}\\
& \textbf{engagiere dich!}\\
& \textbf{engagieren wir uns!}\\
& \textbf{engagiert euch!}\\
& \textbf{engagieren Sie sich!}\\
\end{tabular}

\vspace{10pt}

\begin{tabular}{rl}
\multicolumn{2}{l}{\textbf{PERFEKT}}\\
& haben + sich + \textbf{engagiert}\\
\end{tabular}
\hfill
\begin{tabular}{rl}
\multicolumn{2}{l}{\textbf{FUTURE I}}\\
& werden + sich + \textbf{engagieren}\\
\end{tabular}
\hfill
\begin{tabular}{rl}
\multicolumn{2}{l}{\textbf{PLUSQUAMPERFEKT}}\\
& hatten + sich + \textbf{engagiert}\\
\end{tabular}

\vspace{10pt}

\begin{tabular}{lll}
\multicolumn{3}{l}{\textbf{MIT PRÄPOSITIONEN}}\\
& \textbf{sich engagieren für} + Akk & (to commit oneself to / support)\\
& \textbf{sich engagieren bei} + Dat & (to get involved at / with)\\
\end{tabular}
\hfill
\begin{tabular}{lll}
\multicolumn{3}{l}{\textbf{TRENNBARE VERBEN}}\\
& \textbf{---} & \\
\end{tabular}
\vspace{-5pt}
\end{verbentry}

% sich entscheiden (to decide)
\begin{verbentry}{
  style = {default},
  toc = {sich entscheiden (to decide)},
  name = {sich entscheiden},
  english = {to decide},
  type = {regular},
  auxiliary = {haben}
}
 
    
       \vspace{5pt}

    \hrule
    
    \vspace{5pt}

  \begin{tabular}{rl}
     \multicolumn{2}{l}{\textbf{PRÄSENS} }\\
    ich & \textbf{entscheide mich}\\
    du & \textbf{entscheidest dich}\\
    er/sie/es & \textbf{entscheidet sich}\\
    wir & \textbf{entscheiden uns}\\
    ihr & \textbf{entscheidet euch}\\
    sie/Sie & \textbf{entscheiden sich}\\
  \end{tabular}
  \hfill
     \begin{tabular}{rl}
    \multicolumn{2}{l}{\textbf{PRÄTERITUM} }\\
    ich & \textbf{entschied mich}\\
    du & \textbf{entschiedest dich}\\
    er/sie/es & \textbf{entschied sich}\\
    wir & \textbf{entschieden uns}\\
    ihr & \textbf{entschiedet euch}\\
    sie/Sie & \textbf{entschieden sich}\\
  \end{tabular}
  \hfill
  \begin{tabular}{rl}
     \multicolumn{2}{l}{\textbf{IMPERATIV} }\\
   & \textbf{entscheid(e) (du) dich!}\\
   & \textbf{entscheiden wir uns!}\\
   & \textbf{entscheidet (ihr) euch!}\\
   & \textbf{entscheiden Sie sich!}\\
   & \vspace{10pt}
  \end{tabular}

    \vspace{10pt}

 \begin{tabular}{rl}
     \multicolumn{2}{l}{\textbf{PERFEKT}}\\
    &haben + sich + \textbf{entschieden}\\
  \end{tabular}
\hfill
   \begin{tabular}{rl}
   
    \multicolumn{2}{l}{\textbf{FUTURE I} }\\
    &werden + sich + \textbf{entscheiden}\\
 
  \end{tabular}
\hfill
   \begin{tabular}{rl}
      \multicolumn{2}{l}{\textbf{PLUSQUAMPERFEKT}}\\
    &hatten + sich + \textbf{entschieden}\\
  \end{tabular}
  
\vspace{10pt}

   \begin{tabular}{lll}
      \multicolumn{3}{l}{\textbf{MIT PRÄPOSITIONEN}}\\
& \textbf{sich entscheiden für} + Akk & (to decide for something)\\
& \textbf{sich entscheiden gegen} + Akk & (to decide against something)
  \end{tabular}
  \hfill
   \begin{tabular}{lll}
\multicolumn{3}{l}{\textbf{TRENNBARE VERBEN}}\\
& \textbf{---} & \\
\end{tabular}
  
  \vspace{-5pt}
\end{verbentry}

% sich entschuldigen (to apologise)
\begin{verbentry}{
  style = {acc},
  toc = {sich entschuldigen (to apologise)},
  name = {sich entschuldigen},
  english = {to apologise},
  type = {regular, + Akkusativ},
  auxiliary = {haben}
}
 
    
       \vspace{5pt}

    \hrule
    
    \vspace{5pt}

  \begin{tabular}{rl}
     \multicolumn{2}{l}{\textbf{PRÄSENS} }\\
    ich & \textbf{entschuldige mich}\\
    du & \textbf{entschuldigst dich}\\
    er/sie/es & \textbf{entschuldigt sich}\\
    wir & \textbf{entschuldigen uns}\\
    ihr & \textbf{entschuldigt euch}\\
    sie/Sie & \textbf{entschuldigen sich}\\
  \end{tabular}
  \hfill
     \begin{tabular}{rl}
    \multicolumn{2}{l}{\textbf{PRÄTERITUM} }\\
    ich & \textbf{entschuldigte mich}\\
    du & \textbf{entschuldigtest dich}\\
    er/sie/es & \textbf{entschuldigte sich}\\
    wir & \textbf{entschuldigten uns}\\
    ihr & \textbf{entschuldigtet euch}\\
    sie/Sie & \textbf{entschuldigten sich}\\
  \end{tabular}
  \hfill
  \begin{tabular}{rl}
     \multicolumn{2}{l}{\textbf{IMPERATIV} }\\
   & \textbf{entschuldig(e) (du) dich!}\\
   & \textbf{entschuldigen wir uns!}\\
   & \textbf{entschuldigt (ihr) euch!}\\
   & \textbf{entschuldigen Sie sich!}\\
   & \vspace{10pt}
  \end{tabular}

    \vspace{10pt}

 \begin{tabular}{rl}
     \multicolumn{2}{l}{\textbf{PERFEKT}}\\
    &haben + sich + \textbf{entschuldigt}\\
  \end{tabular}
\hfill
   \begin{tabular}{rl}
   
    \multicolumn{2}{l}{\textbf{FUTURE I} }\\
    &werden + sich + \textbf{entschuldigen}\\
 
  \end{tabular}
\hfill
   \begin{tabular}{rl}
      \multicolumn{2}{l}{\textbf{PLUSQUAMPERFEKT}}\\
    &hatten + sich + \textbf{entschuldigt}\\
  \end{tabular}
  
\vspace{10pt}

   \begin{tabular}{lll}
      \multicolumn{3}{l}{\textbf{MIT PRÄPOSITIONEN}}\\
& \textbf{sich entschuldigen für} + Akk & (to apologise for)\\
& \textbf{sich entschuldigen bei} + Dat & (to apologise to)\\
& \textbf{sich entschuldigen mit} + Dat & (to plead something)
  \end{tabular}
  \hfill
   \begin{tabular}{lll}
      \multicolumn{3}{l}{\textbf{TRENNBARE VERBEN}}\\
\vspace{30pt}
 \end{tabular}
  
  \vspace{-5pt}
\end{verbentry}

% sich entspannen (to relax)
\begin{verbentry}{
  style = {default},
  toc = {sich entspannen (to relax)},
  name = {sich entspannen},
  english = {to relax},
  type = {regular},
  auxiliary = {haben}
}
 
    
     \vspace{5pt}

    \hrule
    
    \vspace{5pt}

  \begin{tabular}{rl}
     \multicolumn{2}{l}{\textbf{PRÄSENS} }\\
    ich & \textbf{entspanne mich}\\
    du & \textbf{entspannst dich}\\
    er/sie/es & \textbf{entspannt sich}\\
    wir & \textbf{entspannen uns}\\
    ihr & \textbf{entspannt euch}\\
    sie/Sie & \textbf{entspannen sich}\\
  \end{tabular}
  \hfill
     \begin{tabular}{rl}
    \multicolumn{2}{l}{\textbf{PRÄTERITUM} }\\
    ich & \textbf{entspannte mich}\\
    du & \textbf{entspanntest dich}\\
    er/sie/es & \textbf{entspannte sich}\\
    wir & \textbf{entspannten uns}\\
    ihr & \textbf{entspanntet euch}\\
    sie/Sie & \textbf{entspannten sich}\\
  \end{tabular}
  \hfill
  \begin{tabular}{rl}
     \multicolumn{2}{l}{\textbf{IMPERATIV} }\\
   & \textbf{entspann dich (du)!}\\
   & \textbf{entspannen wir uns!}\\
   & \textbf{entspannt euch (ihr)!}\\
   & \textbf{entspannen Sie sich!}\\
   & \vspace{10pt}
  \end{tabular}

    \vspace{10pt}

 \begin{tabular}{rl}
     \multicolumn{2}{l}{\textbf{PERFEKT}}\\
   & haben + sich + \textbf{entspannt}\\
  \end{tabular}
\hfill
   \begin{tabular}{rl}
   
    \multicolumn{2}{l}{\textbf{FUTURE I} }\\
   & werden + sich + \textbf{entspannen}\\
 
  \end{tabular}
\hfill
   \begin{tabular}{rl}
      \multicolumn{2}{l}{\textbf{PLUSQUAMPERFEKT}}\\
   & hatten + sich + \textbf{entspannt}\\
  \end{tabular}

   \vspace{10pt}

   \begin{tabular}{lll}
      \multicolumn{3}{l}{\textbf{MIT PRÄPOSITIONEN}}\\
& \textbf{sich entspannen bei} + Dat & (to relax while / during)\\
& \textbf{sich entspannen mit} + Dat & (to relax with)\\
& \textbf{sich entspannen nach} + Dat & (to relax after)\\
\end{tabular}
  \hfill
\begin{tabular}{lll}
\multicolumn{3}{l}{\textbf{TRENNBARE VERBEN}}\\
& \textbf{---} & \\
\end{tabular}


  \vspace{-5pt}
\end{verbentry}

% sich entwickeln (to develop)
\begin{verbentry}{
  style = {default},
  toc = {sich entwickeln (to develop)},
  name = {sich entwickeln},
  english = {to develop},
  type = {regular},
  auxiliary = {haben}
}
 
    
       \vspace{5pt}

    \hrule
    
    \vspace{5pt}

  \begin{tabular}{rl}
     \multicolumn{2}{l}{\textbf{PRÄSENS} }\\
    ich & \textbf{entwickele mich}\\
    du & \textbf{entwickelst dich}\\
    er/sie/es & \textbf{entwickelt sich}\\
    wir & \textbf{entwickeln uns}\\
    ihr & \textbf{entwickelt euch}\\
    sie/Sie & \textbf{entwickeln sich}\\
  \end{tabular}
  \hfill
     \begin{tabular}{rl}
    \multicolumn{2}{l}{\textbf{PRÄTERITUM} }\\
    ich & \textbf{entwickelte mich}\\
    du & \textbf{ventwickeltest dich}\\
    er/sie/es & \textbf{entwickelte sich}\\
    wir & \textbf{entwickelten uns}\\
    ihr & \textbf{entwickeltet euch}\\
    sie/Sie & \textbf{entwickelten sich}\\
  \end{tabular}
  \hfill
  \begin{tabular}{rl}
     \multicolumn{2}{l}{\textbf{IMPERATIV} }\\
   & \textbf{entwickel(e) (du) dich!}\\
   & \textbf{entwickelen wir uns!}\\
   & \textbf{entwickelt (ihr) euch!}\\
   & \textbf{entwickelen Sie sich!}\\
   & \vspace{10pt}
  \end{tabular}

    \vspace{10pt}

 \begin{tabular}{rl}
     \multicolumn{2}{l}{\textbf{PERFEKT}}\\
    &haben + sich + \textbf{entwickelt}\\
  \end{tabular}
\hfill
   \begin{tabular}{rl}
   
    \multicolumn{2}{l}{\textbf{FUTURE I} }\\
    &werden + sich + \textbf{entwickeln}\\
 
  \end{tabular}
\hfill
   \begin{tabular}{rl}
      \multicolumn{2}{l}{\textbf{PLUSQUAMPERFEKT}}\\
    &hatten + sich + \textbf{entwickelt}\\
  \end{tabular}
  
\vspace{10pt}

   \begin{tabular}{lll}
      \multicolumn{3}{l}{\textbf{MIT PRÄPOSITIONEN}}\\
& \textbf{sich entwickeln aus} + Dat & (to evolve from something)\\
& \textbf{sich entwickeln zu} + Dat & (to evolve into something)
  \end{tabular}
  \hfill
   \begin{tabular}{lll}
\multicolumn{3}{l}{\textbf{TRENNBARE VERBEN}}\\
& \textbf{---} & \\
\end{tabular}
  
  \vspace{-5pt}
\end{verbentry}

% sich erholen (to recover / rest)
\begin{verbentry}{
  style = {default},
  toc = {sich erholen (to recover / rest)},
  name = {sich erholen},
  english = {to recover / rest},
  type = {regular},
  auxiliary = {haben}
}


\vspace{5pt}
\hrule
\vspace{5pt}

\begin{tabular}{rl}
\multicolumn{2}{l}{\textbf{PRÄSENS}}\\
ich & \textbf{erhole mich}\\
du & \textbf{erholst dich}\\
er/sie/es & \textbf{erholt sich}\\
wir & \textbf{erholen uns}\\
ihr & \textbf{erholt euch}\\
sie/Sie & \textbf{erholen sich}\\
\end{tabular}
\hfill
\begin{tabular}{rl}
\multicolumn{2}{l}{\textbf{PRÄTERITUM}}\\
ich & \textbf{erholte mich}\\
du & \textbf{erholtest dich}\\
er/sie/es & \textbf{erholte sich}\\
wir & \textbf{erholten uns}\\
ihr & \textbf{erholtet euch}\\
sie/Sie & \textbf{erholten sich}\\
\end{tabular}
\hfill
\begin{tabular}{rl}
\multicolumn{2}{l}{\textbf{IMPERATIV}}\\
& \textbf{erhole dich!}\\
& \textbf{erholen wir uns!}\\
& \textbf{erholt euch!}\\
& \textbf{erholen Sie sich!}\\
\end{tabular}

\vspace{10pt}

\begin{tabular}{rl}
\multicolumn{2}{l}{\textbf{PERFEKT}}\\
& haben + sich + \textbf{erholt}\\
\end{tabular}
\hfill
\begin{tabular}{rl}
\multicolumn{2}{l}{\textbf{FUTURE I}}\\
& werden + sich + \textbf{erholen}\\
\end{tabular}
\hfill
\begin{tabular}{rl}
\multicolumn{2}{l}{\textbf{PLUSQUAMPERFEKT}}\\
& hatten + sich + \textbf{erholt}\\
\end{tabular}

\vspace{10pt}

\begin{tabular}{lll}
\multicolumn{3}{l}{\textbf{MIT PRÄPOSITIONEN}}\\
& \textbf{sich erholen von} + Dat & (to recover from)\\
& \textbf{sich erholen nach} + Dat & (to recover after)\\
\end{tabular}
\hfill
\begin{tabular}{lll}
\multicolumn{3}{l}{\textbf{TRENNBARE VERBEN}}\\
& \textbf{---} & \\
\end{tabular}

\vspace{-5pt}
\end{verbentry}

% sich freuen (to be pleased about)
\begin{verbentry}{
  style = {default},
  toc = {sich freuen (to be pleased about)},
  name = {sich freuen},
  english = {to be pleased about},
  type = {regular},
  auxiliary = {haben}
}
 
    
     \vspace{5pt}

    \hrule
    
    \vspace{5pt}

  \begin{tabular}{rl}
     \multicolumn{2}{l}{\textbf{PRÄSENS} }\\
    ich & \textbf{freue mich}\\
    du & \textbf{freust dich}\\
    er/sie/es & \textbf{freut sich}\\
    wir & \textbf{freuen uns}\\
    ihr & \textbf{freut euch}\\
    sie/Sie & \textbf{freuen sich}\\
  \end{tabular}
  \hfill
     \begin{tabular}{rl}
    \multicolumn{2}{l}{\textbf{PRÄTERITUM} }\\
    ich & \textbf{freute mich}\\
    du & \textbf{freutest dich}\\
    er/sie/es & \textbf{freute sich}\\
    wir & \textbf{freuten uns}\\
    ihr & \textbf{freutet euch}\\
    sie/Sie & \textbf{freiten sich}\\
  \end{tabular}
  \hfill
  \begin{tabular}{rl}
     \multicolumn{2}{l}{\textbf{IMPERATIV} }\\
   & \textbf{freu(e) (du) dich!}\\
   & \textbf{freuen wir uns!}\\
   & \textbf{freut (ihr) euch!}\\
   & \textbf{freuen Sie sich!}\\
   & \vspace{10pt}
  \end{tabular}

    \vspace{10pt}

 \begin{tabular}{rl}
     \multicolumn{2}{l}{\textbf{PERFEKT}}\\
   &  haben + sich + \textbf{gefreut}\\
  \end{tabular}
\hfill
   \begin{tabular}{rl}
    \multicolumn{2}{l}{\textbf{FUTURE I} }\\
   & werden + sich + \textbf{freuen}\\
  \end{tabular}
\hfill
   \begin{tabular}{rl}
      \multicolumn{2}{l}{\textbf{PLUSQUAMPERFEKT}}\\
    & hatten + sich + \textbf{gefreut}\\
  \end{tabular}
  
\vspace{10pt}

   \begin{tabular}{lll}
      \multicolumn{3}{l}{\textbf{MIT PRÄPOSITIONEN}}\\
& \textbf{sich freuen über} + Akk & (to be please about something)\\
& \textbf{sich freuen auf} + Akk & (to look forward to something)
  \end{tabular}
  \hfill
   \begin{tabular}{lll}
\multicolumn{3}{l}{\textbf{TRENNBARE VERBEN}}\\
& \textbf{---} & \\
\end{tabular}



  \vspace{-5pt}
\end{verbentry}

% sich fühlen (to feel)
\begin{verbentry}{
  style = {acc},
  toc = {sich fühlen (to feel)},
  name = {sich fühlen},
  english = {to feel},
  type = {regular, + Akkusativ},
  auxiliary = {haben}
}
 
    
       \vspace{5pt}

    \hrule
    
    \vspace{5pt}

  \begin{tabular}{rl}
     \multicolumn{2}{l}{\textbf{PRÄSENS} }\\
    ich & \textbf{fühle mich}\\
    du & \textbf{fühlst dich}\\
    er/sie/es & \textbf{fühlt sich}\\
    wir & \textbf{fühlen uns}\\
    ihr & \textbf{fühlt euch}\\
    sie/Sie & \textbf{fühlen sich}\\
  \end{tabular}
  \hfill
     \begin{tabular}{rl}
    \multicolumn{2}{l}{\textbf{PRÄTERITUM} }\\
    ich & \textbf{fühlte mich}\\
    du & \textbf{fühltest dich}\\
    er/sie/es & \textbf{fühlte sich}\\
    wir & \textbf{fühlten uns}\\
    ihr & \textbf{fühltet euch}\\
    sie/Sie & \textbf{fühlten sich}\\
  \end{tabular}
  \hfill
  \begin{tabular}{rl}
     \multicolumn{2}{l}{\textbf{IMPERATIV} }\\
   & \textbf{fühl(e) (du) dich!}\\
   & \textbf{fühlen wir uns!}\\
   & \textbf{fühlt (ihr) euch!}\\
   & \textbf{fühlen Sie sich!}\\
   & \vspace{10pt}
  \end{tabular}

    \vspace{10pt}

 \begin{tabular}{rl}
     \multicolumn{2}{l}{\textbf{PERFEKT}}\\
    &haben + sich + \textbf{gefühlt}\\
  \end{tabular}
\hfill
   \begin{tabular}{rl}
   
    \multicolumn{2}{l}{\textbf{FUTURE I} }\\
    &werden + sich + \textbf{fühlen}\\
 
  \end{tabular}
\hfill
   \begin{tabular}{rl}
      \multicolumn{2}{l}{\textbf{PLUSQUAMPERFEKT}}\\
    &hatten + sich + \textbf{gefühlt}\\
  \end{tabular}
  
\vspace{10pt}

\begin{tabular}{lll}
\multicolumn{3}{l}{\textbf{MIT PRÄPOSITIONEN}}\\
& \textbf{---} & \\
\end{tabular}
  \hfill
   \begin{tabular}{lll}
      \multicolumn{3}{l}{\textbf{TRENNBARE VERBEN}}\\
     & \textbf{sich wohl$|$fühlen} & (to be comfortable)
  \end{tabular}
  
  \vspace{-5pt}
\end{verbentry}

% sich fürchten (to be afraid)
\begin{verbentry}{
  style = {default},
  toc = {sich fürchten (to be afraid)},
  name = {sich fürchten},
  english = {to be afraid},
  type = {regular},
  auxiliary = {haben}
}
 
    
     \vspace{5pt}

    \hrule
    
    \vspace{5pt}

  \begin{tabular}{rl}
     \multicolumn{2}{l}{\textbf{PRÄSENS} }\\
    ich & \textbf{fürchte mich}\\
    du & \textbf{fürchtest dich}\\
    er/sie/es & \textbf{fürchtet sich}\\
    wir & \textbf{fürchten uns}\\
    ihr & \textbf{fürchtet euch}\\
    sie/Sie & \textbf{fürchten sich}\\
  \end{tabular}
  \hfill
     \begin{tabular}{rl}
    \multicolumn{2}{l}{\textbf{PRÄTERITUM} }\\
    ich & \textbf{fürchtete mich}\\
    du & \textbf{fürchtetest dich}\\
    er/sie/es & \textbf{fürchtete sich}\\
    wir & \textbf{fürchteten uns}\\
    ihr & \textbf{fürchtetet euch}\\
    sie/Sie & \textbf{fürchteten sich}\\
  \end{tabular}
  \hfill
  \begin{tabular}{rl}
     \multicolumn{2}{l}{\textbf{IMPERATIV} }\\
   & \textbf{fürchte dich!}\\
   & \textbf{fürchten wir uns!}\\
   & \textbf{fürchtet euch!}\\
   & \textbf{fürchten Sie sich!}\\
   & \vspace{10pt}
  \end{tabular}

    \vspace{10pt}

 \begin{tabular}{rl}
     \multicolumn{2}{l}{\textbf{PERFEKT}}\\
   & haben + sich + \textbf{gefürchtet}\\
  \end{tabular}
\hfill
   \begin{tabular}{rl}
    \multicolumn{2}{l}{\textbf{FUTURE I} }\\
   & werden + sich + \textbf{fürchten}\\
  \end{tabular}
\hfill
   \begin{tabular}{rl}
      \multicolumn{2}{l}{\textbf{PLUSQUAMPERFEKT}}\\
   & hatten + sich + \textbf{gefürchtet}\\
  \end{tabular}

   \vspace{10pt}

   \begin{tabular}{lll}
      \multicolumn{3}{l}{\textbf{MIT PRÄPOSITIONEN}}\\
& \textbf{sich fürchten vor} + Dat & (to be afraid of)\\
& \textbf{sich fürchten um} + Akk & (to fear for)\\
\end{tabular}
  \hfill
   \begin{tabular}{lll}
\multicolumn{3}{l}{\textbf{TRENNBARE VERBEN}}\\
& \textbf{---} & \\
\end{tabular}

  \vspace{-5pt}
\end{verbentry}

% sich gewöhnen (to get used to)
\begin{verbentry}{
  style = {default},
  toc = {sich gewöhnen (to get used to)},
  name = {sich gewöhnen},
  english = {to get used to},
  type = {regular},
  auxiliary = {haben}
}
 
    
     \vspace{5pt}

    \hrule
    
    \vspace{5pt}

  \begin{tabular}{rl}
     \multicolumn{2}{l}{\textbf{PRÄSENS} }\\
    ich & \textbf{gewöhne mich}\\
    du & \textbf{gewöhnst dich}\\
    er/sie/es & \textbf{gewöhnt sich}\\
    wir & \textbf{gewöhnen uns}\\
    ihr & \textbf{gewöhnt euch}\\
    sie/Sie & \textbf{gewöhnen sich}\\
  \end{tabular}
  \hfill
     \begin{tabular}{rl}
    \multicolumn{2}{l}{\textbf{PRÄTERITUM} }\\
    ich & \textbf{gewöhnte mich}\\
    du & \textbf{gewöhntest dich}\\
    er/sie/es & \textbf{gewöhnte sich}\\
    wir & \textbf{gewöhnten uns}\\
    ihr & \textbf{gewöhntet euch}\\
    sie/Sie & \textbf{gewöhnten sich}\\
  \end{tabular}
  \hfill
  \begin{tabular}{rl}
     \multicolumn{2}{l}{\textbf{IMPERATIV} }\\
   & \textbf{gewöhn(e) dich!}\\
   & \textbf{gewöhnen wir uns!}\\
   & \textbf{gewöhnt euch!}\\
   & \textbf{gewöhnen Sie sich!}\\
   & \vspace{10pt}
  \end{tabular}

    \vspace{10pt}

 \begin{tabular}{rl}
     \multicolumn{2}{l}{\textbf{PERFEKT}}\\
   & haben + sich + \textbf{gewöhnt}\\
  \end{tabular}
\hfill
   \begin{tabular}{rl}
    \multicolumn{2}{l}{\textbf{FUTURE I} }\\
   & werden + sich + \textbf{gewöhnen}\\
  \end{tabular}
\hfill
   \begin{tabular}{rl}
      \multicolumn{2}{l}{\textbf{PLUSQUAMPERFEKT}}\\
   & hatten + sich + \textbf{gewöhnt}\\
  \end{tabular}

   \vspace{10pt}

   \begin{tabular}{lll}
      \multicolumn{3}{l}{\textbf{MIT PRÄPOSITIONEN}}\\
& \textbf{sich gewöhnen an} + Akk & (to get used to)\\
\end{tabular}
  \hfill
   \begin{tabular}{lll}
\multicolumn{3}{l}{\textbf{TRENNBARE VERBEN}}\\
& \textbf{---} & \\
\end{tabular}

  \vspace{-5pt}
\end{verbentry}

% sich hüten (to beware / to be careful)
\begin{verbentry}{
  style = {default},
  toc = {sich hüten (to beware / to be careful)},
  name = {sich hüten},
  english = {to beware / to be careful},
  type = {regular},
  auxiliary = {haben}
}
 
    
     \vspace{5pt}

    \hrule
    
    \vspace{5pt}

  \begin{tabular}{rl}
     \multicolumn{2}{l}{\textbf{PRÄSENS} }\\
    ich & \textbf{hüte mich}\\
    du & \textbf{hütest dich}\\
    er/sie/es & \textbf{hütet sich}\\
    wir & \textbf{hüten uns}\\
    ihr & \textbf{hütet euch}\\
    sie/Sie & \textbf{hüten sich}\\
  \end{tabular}
  \hfill
     \begin{tabular}{rl}
    \multicolumn{2}{l}{\textbf{PRÄTERITUM} }\\
    ich & \textbf{hütete mich}\\
    du & \textbf{hütetest dich}\\
    er/sie/es & \textbf{hütete sich}\\
    wir & \textbf{hüteten uns}\\
    ihr & \textbf{hütetet euch}\\
    sie/Sie & \textbf{hüteten sich}\\
  \end{tabular}
  \hfill
  \begin{tabular}{rl}
     \multicolumn{2}{l}{\textbf{IMPERATIV} }\\
   & \textbf{hüte dich!}\\
   & \textbf{hüten wir uns!}\\
   & \textbf{hütet euch!}\\
   & \textbf{hüten Sie sich!}\\
   & \vspace{10pt}
  \end{tabular}

    \vspace{10pt}

 \begin{tabular}{rl}
     \multicolumn{2}{l}{\textbf{PERFEKT}}\\
   & haben + sich + \textbf{gehütet}\\
  \end{tabular}
\hfill
   \begin{tabular}{rl}
    \multicolumn{2}{l}{\textbf{FUTURE I} }\\
   & werden + sich + \textbf{hüten}\\
  \end{tabular}
\hfill
   \begin{tabular}{rl}
      \multicolumn{2}{l}{\textbf{PLUSQUAMPERFEKT}}\\
   & hatten + sich + \textbf{gehütet}\\
  \end{tabular}

   \vspace{10pt}

   \begin{tabular}{lll}
      \multicolumn{3}{l}{\textbf{MIT PRÄPOSITIONEN}}\\
& \textbf{sich hüten vor} + Dat & (to beware of)\\
& \textbf{sich hüten, zu} + Inf & (to be careful not to)\\
\end{tabular}
  \hfill
   \begin{tabular}{lll}
\multicolumn{3}{l}{\textbf{TRENNBARE VERBEN}}\\
& \textbf{---} & \\
\end{tabular}

  \vspace{-5pt}
\end{verbentry}

% sich interessieren (to be interested in)
\begin{verbentry}{
  style = {default},
  toc = {sich interessieren (to be interested in)},
  name = {sich interessieren},
  english = {to be interest in},
  type = {regular},
  auxiliary = {haben}
}
 
    
     \vspace{5pt}

    \hrule
    
    \vspace{5pt}

  \begin{tabular}{rl}
     \multicolumn{2}{l}{\textbf{PRÄSENS} }\\
    ich & \textbf{interessiere mich}\\
    du & \textbf{interessierst dich}\\
    er/sie/es & \textbf{interessiert sich}\\
    wir & \textbf{interessieren uns}\\
    ihr & \textbf{interessiert euch}\\
    sie/Sie & \textbf{interessieren sich}\\
  \end{tabular}
  \hfill
     \begin{tabular}{rl}
    \multicolumn{2}{l}{\textbf{PRÄTERITUM} }\\
    ich & \textbf{interessierte mich}\\
    du & \textbf{interessiertest dich}\\
    er/sie/es & \textbf{interessierte sich}\\
    wir & \textbf{interessierten uns}\\
    ihr & \textbf{interessiertet euch}\\
    sie/Sie & \textbf{interessierten sich}\\
  \end{tabular}
  \hfill
  \begin{tabular}{rl}
     \multicolumn{2}{l}{\textbf{IMPERATIV} }\\
   & \textbf{interessier(e) (du) dich!}\\
   & \textbf{interessieren wir uns!}\\
   & \textbf{interessiert (ihr) euch!}\\
   & \textbf{interessieren Sie sich!}\\
   & \vspace{10pt}
  \end{tabular}

    \vspace{10pt}

 \begin{tabular}{rl}
     \multicolumn{2}{l}{\textbf{PERFEKT}}\\
   & haben + sich + \textbf{interessiert}\\
  \end{tabular}
\hfill
   \begin{tabular}{rl}
    \multicolumn{2}{l}{\textbf{FUTURE I} }\\
    &werden + sich + \textbf{interessieren}\\
  \end{tabular}
\hfill
   \begin{tabular}{rl}
      \multicolumn{2}{l}{\textbf{PLUSQUAMPERFEKT}}\\
   & hatten + sich + \textbf{interessiert}\\
  \end{tabular}
  
\vspace{10pt}

   \begin{tabular}{lll}
      \multicolumn{3}{l}{\textbf{MIT PRÄPOSITIONEN}}\\
& \textbf{sich interessieren für} + Akk & (to be interested in something)\\
  \end{tabular}
\hfill
\begin{tabular}{lll}
\multicolumn{3}{l}{\textbf{TRENNBARE VERBEN}}\\
& \textbf{---} & \\
\end{tabular}
 % \hfill
 %  \begin{tabular}{lll}
  %    \multicolumn{3}{l}{\textbf{TRENNBARE VERBEN}}\\
%\vspace{10pt}
 % \end{tabular}

  \vspace{-5pt}
\end{verbentry}

% sich irren (to be mistaken)
\begin{verbentry}{
  style = {default},
  toc = {sich irren (to be mistaken)},
  name = {sich irren},
  english = {to be mistaken},
  type = {regular},
  auxiliary = {haben}
}
 
    
     \vspace{5pt}

    \hrule
    
    \vspace{5pt}

  \begin{tabular}{rl}
     \multicolumn{2}{l}{\textbf{PRÄSENS} }\\
    ich & \textbf{irre mich}\\
    du & \textbf{irrst dich}\\
    er/sie/es & \textbf{irrt sich}\\
    wir & \textbf{irren uns}\\
    ihr & \textbf{irrt euch}\\
    sie/Sie & \textbf{irren sich}\\
  \end{tabular}
  \hfill
     \begin{tabular}{rl}
    \multicolumn{2}{l}{\textbf{PRÄTERITUM} }\\
    ich & \textbf{irrte mich}\\
    du & \textbf{irrtest dich}\\
    er/sie/es & \textbf{irrte sich}\\
    wir & \textbf{irrten uns}\\
    ihr & \textbf{irrtet euch}\\
    sie/Sie & \textbf{irrten sich}\\
  \end{tabular}
  \hfill
  \begin{tabular}{rl}
     \multicolumn{2}{l}{\textbf{IMPERATIV} }\\
   & \textbf{irr(e) dich!}\\
   & \textbf{irren wir uns!}\\
   & \textbf{irrt euch!}\\
   & \textbf{irren Sie sich!}\\
   & \vspace{10pt}
  \end{tabular}

    \vspace{10pt}

 \begin{tabular}{rl}
     \multicolumn{2}{l}{\textbf{PERFEKT}}\\
   & haben + sich + \textbf{geirrt}\\
  \end{tabular}
\hfill
   \begin{tabular}{rl}
    \multicolumn{2}{l}{\textbf{FUTURE I} }\\
   & werden + sich + \textbf{irren}\\
  \end{tabular}
\hfill
   \begin{tabular}{rl}
      \multicolumn{2}{l}{\textbf{PLUSQUAMPERFEKT}}\\
   & hatten + sich + \textbf{geirrt}\\
  \end{tabular}

   \vspace{10pt}

   \begin{tabular}{lll}
      \multicolumn{3}{l}{\textbf{MIT PRÄPOSITIONEN}}\\
& \textbf{sich irren in} + Dat & (to be mistaken about)\\
& \textbf{sich irren bei} + Dat & (to make a mistake in/with)\\
\end{tabular}
  \hfill
   \begin{tabular}{lll}
\multicolumn{3}{l}{\textbf{TRENNBARE VERBEN}}\\
& \textbf{---} & \\
\end{tabular}

  \vspace{-5pt}
\end{verbentry}

% sich konzentrieren (to concentrate)
\begin{verbentry}{
  style = {default},
  toc = {sich konzentrieren (to concentrate)},
  name = {sich konzentrieren},
  english = {to concentrate},
  type = {regular},
  auxiliary = {haben}
}
 
    
       \vspace{5pt}

    \hrule
    
    \vspace{5pt}

  \begin{tabular}{rl}
     \multicolumn{2}{l}{\textbf{PRÄSENS} }\\
    ich & \textbf{konzentriere mich}\\
    du & \textbf{konzentrierst dich}\\
    er/sie/es & \textbf{konzentriert sich}\\
    wir & \textbf{konzentrieren uns}\\
    ihr & \textbf{konzentriert euch}\\
    sie/Sie & \textbf{konzentrieren sich}\\
  \end{tabular}
  \hfill
     \begin{tabular}{rl}
    \multicolumn{2}{l}{\textbf{PRÄTERITUM} }\\
    ich & \textbf{konzentrierte mich}\\
    du & \textbf{konzentriertest dich}\\
    er/sie/es & \textbf{konzentrierte sich}\\
    wir & \textbf{konzentrierten uns}\\
    ihr & \textbf{konzentriertet euch}\\
    sie/Sie & \textbf{konzentrierten sich}\\
  \end{tabular}
  \hfill
  \begin{tabular}{rl}
     \multicolumn{2}{l}{\textbf{IMPERATIV} }\\
   & \textbf{konzentrier(e) (du) dich!}\\
   & \textbf{konzentrieren wir uns!}\\
   & \textbf{konzentriert (ihr) euch!}\\
   & \textbf{konzentrieren Sie sich!}\\
   & \vspace{10pt}
  \end{tabular}

    \vspace{10pt}

 \begin{tabular}{rl}
     \multicolumn{2}{l}{\textbf{PERFEKT}}\\
    &haben + sich + \textbf{konzentriert}\\
  \end{tabular}
\hfill
   \begin{tabular}{rl}
   
    \multicolumn{2}{l}{\textbf{FUTURE I} }\\
    &werden + sich + \textbf{konzentrieren}\\
 
  \end{tabular}
\hfill
   \begin{tabular}{rl}
      \multicolumn{2}{l}{\textbf{PLUSQUAMPERFEKT}}\\
    &hatten + sich + \textbf{konzentriert}\\
  \end{tabular}
  
\vspace{10pt}

   \begin{tabular}{lll}
      \multicolumn{3}{l}{\textbf{MIT PRÄPOSITIONEN}}\\
& \textbf{sich konzentrieren auf} + Akk & (to concentrate on something)
  \end{tabular}
  \hfill
   \begin{tabular}{lll}
\multicolumn{3}{l}{\textbf{TRENNBARE VERBEN}}\\
& \textbf{---} & \\
\end{tabular}
  
  \vspace{-5pt}
\end{verbentry}

% sich kümmern (to take care of)
\begin{verbentry}{
  style = {default},
  toc = {sich kümmern (to take care of)},
  name = {sich kümmern},
  english = {to take care of},
  type = {regular},
  auxiliary = {haben}
}
 
    
       \vspace{5pt}

    \hrule
    
    \vspace{5pt}

  \begin{tabular}{rl}
     \multicolumn{2}{l}{\textbf{PRÄSENS} }\\
    ich & \textbf{kümmere mich}\\
    du & \textbf{kümmerst dich}\\
    er/sie/es & \textbf{kümmert sich}\\
    wir & \textbf{kümmern uns}\\
    ihr & \textbf{kümmert euch}\\
    sie/Sie & \textbf{kümmern sich}\\
  \end{tabular}
  \hfill
     \begin{tabular}{rl}
    \multicolumn{2}{l}{\textbf{PRÄTERITUM} }\\
    ich & \textbf{kümmerte mich}\\
    du & \textbf{kümmertest dich}\\
    er/sie/es & \textbf{kümmerte sich}\\
    wir & \textbf{kümmerten uns}\\
    ihr & \textbf{kümmertet euch}\\
    sie/Sie & \textbf{kümmerten sich}\\
  \end{tabular}
  \hfill
  \begin{tabular}{rl}
     \multicolumn{2}{l}{\textbf{IMPERATIV} }\\
   & \textbf{kümmer(e) (du) dich!}\\
   & \textbf{kümmern wir uns!}\\
   & \textbf{kümmert (ihr) euch!}\\
   & \textbf{kümmern Sie sich!}\\
   & \vspace{10pt}
  \end{tabular}

    \vspace{10pt}

 \begin{tabular}{rl}
     \multicolumn{2}{l}{\textbf{PERFEKT}}\\
    &haben + sich + \textbf{gekümmert}\\
  \end{tabular}
\hfill
   \begin{tabular}{rl}
   
    \multicolumn{2}{l}{\textbf{FUTURE I} }\\
    &werden + sich + \textbf{kümmern}\\
 
  \end{tabular}
\hfill
   \begin{tabular}{rl}
      \multicolumn{2}{l}{\textbf{PLUSQUAMPERFEKT}}\\
    &hatten + sich + \textbf{gekümmert}\\
  \end{tabular}
  
\vspace{10pt}

   \begin{tabular}{lll}
      \multicolumn{3}{l}{\textbf{MIT PRÄPOSITIONEN}}\\
& \textbf{sich kümmern um} + Akk & (to take care of)
  \end{tabular}
  \hfill
   \begin{tabular}{lll}
\multicolumn{3}{l}{\textbf{TRENNBARE VERBEN}}\\
& \textbf{---} & \\
\end{tabular}
  
  \vspace{-5pt}
\end{verbentry}

% sich melden (to get in touch / report)
\begin{verbentry}{
  style = {default},
  toc = {sich melden (to get in touch / report)},
  name = {sich melden},
  english = {to get in touch},
  type = {regular},
  auxiliary = {haben}
}


\vspace{5pt}
\hrule
\vspace{5pt}

\begin{tabular}{rl}
\multicolumn{2}{l}{\textbf{PRÄSENS}}\\
ich & \textbf{melde mich}\\
du & \textbf{meldest dich}\\
er/sie/es & \textbf{meldet sich}\\
wir & \textbf{melden uns}\\
ihr & \textbf{meldet euch}\\
sie/Sie & \textbf{melden sich}\\
\end{tabular}
\hfill
\begin{tabular}{rl}
\multicolumn{2}{l}{\textbf{PRÄTERITUM}}\\
ich & \textbf{meldete mich}\\
du & \textbf{meldetest dich}\\
er/sie/es & \textbf{meldete sich}\\
wir & \textbf{meldeten uns}\\
ihr & \textbf{meldetet euch}\\
sie/Sie & \textbf{meldeten sich}\\
\end{tabular}
\hfill
\begin{tabular}{rl}
\multicolumn{2}{l}{\textbf{IMPERATIV}}\\
& \textbf{melde dich (du)!}\\
& \textbf{melden wir uns!}\\
& \textbf{meldet euch (ihr)!}\\
& \textbf{melden Sie sich!}\\
\end{tabular}

\vspace{10pt}

\begin{tabular}{rl}
\multicolumn{2}{l}{\textbf{PERFEKT}}\\
& haben + sich + \textbf{gemeldet}\\
\end{tabular}
\hfill
\begin{tabular}{rl}
\multicolumn{2}{l}{\textbf{FUTURE I}}\\
& werden + sich +\textbf{melden}\\
\end{tabular}
\hfill
\begin{tabular}{rl}
\multicolumn{2}{l}{\textbf{PLUSQUAMPERFEKT}}\\
& hatten + sich + \textbf{gemeldet}\\
\end{tabular}

\vspace{10pt}

\begin{tabular}{lll}
\multicolumn{3}{l}{\textbf{MIT PRÄPOSITIONEN}}\\
& \textbf{sich melden bei} + Dat & (to contact)\\
& \textbf{sich melden für} + Akk & (to register for)\\
\end{tabular}
\hfill
\begin{tabular}{lll}
\multicolumn{3}{l}{\textbf{TRENNBARE VERBEN}}\\
& \textbf{an$|$melden} & (to register)\\
& \textbf{ab$|$melden} & (to deregister)\\
\end{tabular}
\vspace{-5pt}
\end{verbentry}

% sich rächen (to take revenge)
\begin{verbentry}{
  style = {default},
  toc = {sich rächen (to take revenge)},
  name = {sich rächen},
  english = {to take revenge},
  type = {regular},
  auxiliary = {haben}
}
 
    
     \vspace{5pt}

    \hrule
    
    \vspace{5pt}

  \begin{tabular}{rl}
     \multicolumn{2}{l}{\textbf{PRÄSENS} }\\
    ich & \textbf{räche mich}\\
    du & \textbf{rächst dich}\\
    er/sie/es & \textbf{rächt sich}\\
    wir & \textbf{rächen uns}\\
    ihr & \textbf{rächt euch}\\
    sie/Sie & \textbf{rächen sich}\\
  \end{tabular}
  \hfill
     \begin{tabular}{rl}
    \multicolumn{2}{l}{\textbf{PRÄTERITUM} }\\
    ich & \textbf{rächte mich}\\
    du & \textbf{rächtest dich}\\
    er/sie/es & \textbf{rächte sich}\\
    wir & \textbf{rächten uns}\\
    ihr & \textbf{rächtet euch}\\
    sie/Sie & \textbf{rächten sich}\\
  \end{tabular}
  \hfill
  \begin{tabular}{rl}
     \multicolumn{2}{l}{\textbf{IMPERATIV} }\\
   & \textbf{räch(e) dich!}\\
   & \textbf{rächen wir uns!}\\
   & \textbf{rächt euch!}\\
   & \textbf{rächen Sie sich!}\\
   & \vspace{10pt}
  \end{tabular}

    \vspace{10pt}

 \begin{tabular}{rl}
     \multicolumn{2}{l}{\textbf{PERFEKT}}\\
   & haben + sich + \textbf{gerächt}\\
  \end{tabular}
\hfill
   \begin{tabular}{rl}
    \multicolumn{2}{l}{\textbf{FUTURE I} }\\
   & werden + sich + \textbf{rächen}\\
  \end{tabular}
\hfill
   \begin{tabular}{rl}
      \multicolumn{2}{l}{\textbf{PLUSQUAMPERFEKT}}\\
   & hatten + sich + \textbf{gerächt}\\
  \end{tabular}

   \vspace{10pt}

   \begin{tabular}{lll}
      \multicolumn{3}{l}{\textbf{MIT PRÄPOSITIONEN}}\\
& \textbf{sich rächen an} + Dat & (to take revenge on)\\
& \textbf{sich rächen für} + Akk & (to take revenge for)\\
\end{tabular}
  \hfill
   \begin{tabular}{lll}
\multicolumn{3}{l}{\textbf{TRENNBARE VERBEN}}\\
& \textbf{---} & \\
\end{tabular}

  \vspace{-5pt}
\end{verbentry}

% sich regen (to stir / to move)
\begin{verbentry}{
  style = {default},
  toc = {sich regen (to stir / to move)},
  name = {sich regen},
  english = {to stir / to move},
  type = {regular},
  auxiliary = {haben}
}
 
    
     \vspace{5pt}

    \hrule
    
    \vspace{5pt}

  \begin{tabular}{rl}
     \multicolumn{2}{l}{\textbf{PRÄSENS} }\\
    ich & \textbf{rege mich}\\
    du & \textbf{regst dich}\\
    er/sie/es & \textbf{regt sich}\\
    wir & \textbf{regen uns}\\
    ihr & \textbf{regt euch}\\
    sie/Sie & \textbf{regen sich}\\
  \end{tabular}
  \hfill
     \begin{tabular}{rl}
    \multicolumn{2}{l}{\textbf{PRÄTERITUM} }\\
    ich & \textbf{regte mich}\\
    du & \textbf{regtest dich}\\
    er/sie/es & \textbf{regte sich}\\
    wir & \textbf{regten uns}\\
    ihr & \textbf{regtet euch}\\
    sie/Sie & \textbf{regten sich}\\
  \end{tabular}
  \hfill
  \begin{tabular}{rl}
     \multicolumn{2}{l}{\textbf{IMPERATIV} }\\
   & \textbf{reg(e) dich!}\\
   & \textbf{regen wir uns!}\\
   & \textbf{regt euch!}\\
   & \textbf{regen Sie sich!}\\
   & \vspace{10pt}
  \end{tabular}

    \vspace{10pt}

 \begin{tabular}{rl}
     \multicolumn{2}{l}{\textbf{PERFEKT}}\\
   & haben + sich + \textbf{geregt}\\
  \end{tabular}
\hfill
   \begin{tabular}{rl}
    \multicolumn{2}{l}{\textbf{FUTURE I} }\\
   & werden + sich + \textbf{regen}\\
  \end{tabular}
\hfill
   \begin{tabular}{rl}
      \multicolumn{2}{l}{\textbf{PLUSQUAMPERFEKT}}\\
   & hatten + sich + \textbf{geregt}\\
  \end{tabular}

   \vspace{10pt}

   \begin{tabular}{lll}
      \multicolumn{3}{l}{\textbf{MIT PRÄPOSITIONEN}}\\
& \textbf{sich regen in} + Dat & (to stir in / move in)\\
& \textbf{sich regen bei} + Dat & (to stir at / during)\\
\end{tabular}
  \hfill
   \begin{tabular}{lll}
      \multicolumn{3}{l}{\textbf{TRENNBARE VERBEN}}\\
 & \textbf{an$|$regen} & (to stimulate / suggest)\\
 & \textbf{auf$|$regen} & (to upset / excite)\\
\vspace{10pt}
  \end{tabular}

  \vspace{-5pt}
\end{verbentry}

% sich schämen (to be ashamed)
\begin{verbentry}{
  style = {acc},
  toc = {sich schämen (to be ashamed)},
  name = {sich schämen},
  english = {to be ashamed},
  type = {regular, + Akkusativ},
  auxiliary = {haben}
}
 
    
       \vspace{5pt}

    \hrule
    
    \vspace{5pt}

  \begin{tabular}{rl}
     \multicolumn{2}{l}{\textbf{PRÄSENS} }\\
    ich & \textbf{schäme mich}\\
    du & \textbf{schämst dich}\\
    er/sie/es & \textbf{schämt sich}\\
    wir & \textbf{schämen uns}\\
    ihr & \textbf{schämt euch}\\
    sie/Sie & \textbf{schämen sich}\\
  \end{tabular}
  \hfill
     \begin{tabular}{rl}
    \multicolumn{2}{l}{\textbf{PRÄTERITUM} }\\
    ich & \textbf{schämte mich}\\
    du & \textbf{schämtest dich}\\
    er/sie/es & \textbf{schämte sich}\\
    wir & \textbf{schämten uns}\\
    ihr & \textbf{schämtet euch}\\
    sie/Sie & \textbf{schämten sich}\\
  \end{tabular}
  \hfill
  \begin{tabular}{rl}
     \multicolumn{2}{l}{\textbf{IMPERATIV} }\\
   & \textbf{schäm(e) (du) dich!}\\
   & \textbf{schämen wir uns!}\\
   & \textbf{schämt (ihr) euch!}\\
   & \textbf{schämen Sie sich!}\\
   & \vspace{10pt}
  \end{tabular}

    \vspace{10pt}

 \begin{tabular}{rl}
     \multicolumn{2}{l}{\textbf{PERFEKT}}\\
    &haben + sich + \textbf{geschämt}\\
  \end{tabular}
\hfill
   \begin{tabular}{rl}
   
    \multicolumn{2}{l}{\textbf{FUTURE I} }\\
    &werden + sich + \textbf{schämen}\\
 
  \end{tabular}
\hfill
   \begin{tabular}{rl}
      \multicolumn{2}{l}{\textbf{PLUSQUAMPERFEKT}}\\
    &hatten + sich + \textbf{geschämt}\\
  \end{tabular}
  
\vspace{10pt}

   \begin{tabular}{lll}
      \multicolumn{3}{l}{\textbf{MIT PRÄPOSITIONEN}}\\
& \textbf{sich schämen für} + Akk & (to be ashamed of)
  \end{tabular}
\hfill
\begin{tabular}{lll}
\multicolumn{3}{l}{\textbf{TRENNBARE VERBEN}}\\
& \textbf{---} & \\
\end{tabular}
 % \hfill
 %  \begin{tabular}{lll}
  %    \multicolumn{3}{l}{\textbf{TRENNBARE VERBEN}}\\
%\vspace{10pt}
 % \end{tabular}
  
  \vspace{-5pt}
\end{verbentry}

% sich schminken (to put on makeup)
\begin{verbentry}{
  style = {default},
  toc = {sich schminken (to put on makeup)},
  name = {sich schminken},
  english = {to put on makeup},
  type = {regular},
  auxiliary = {haben}
}


\vspace{5pt}
\hrule
\vspace{5pt}

\begin{tabular}{rl}
\multicolumn{2}{l}{\textbf{PRÄSENS}}\\
ich & \textbf{schminke mich}\\
du & \textbf{schminkst dich}\\
er/sie/es & \textbf{schminkt sich}\\
wir & \textbf{schminken uns}\\
ihr & \textbf{schminkt euch}\\
sie/Sie & \textbf{schminken sich}\\
\end{tabular}
\hfill
\begin{tabular}{rl}
\multicolumn{2}{l}{\textbf{PRÄTERITUM}}\\
ich & \textbf{schminkte mich}\\
du & \textbf{schminktest dich}\\
er/sie/es & \textbf{schminkte sich}\\
wir & \textbf{schminkten uns}\\
ihr & \textbf{schminktet euch}\\
sie/Sie & \textbf{schminkten sich}\\
\end{tabular}
\hfill
\begin{tabular}{rl}
\multicolumn{2}{l}{\textbf{IMPERATIV}}\\
& \textbf{schminke dich!}\\
& \textbf{schminken wir uns!}\\
& \textbf{schminkt euch!}\\
& \textbf{schminken Sie sich!}\\
\end{tabular}

\vspace{10pt}

\begin{tabular}{rl}
\multicolumn{2}{l}{\textbf{PERFEKT}}\\
& haben + sich + \textbf{geschminkt}\\
\end{tabular}
\hfill
\begin{tabular}{rl}
\multicolumn{2}{l}{\textbf{FUTURE I}}\\
& werden + sich + \textbf{schminken}\\
\end{tabular}
\hfill
\begin{tabular}{rl}
\multicolumn{2}{l}{\textbf{PLUSQUAMPERFEKT}}\\
& hatten + sich + \textbf{geschminkt}\\
\end{tabular}

\vspace{10pt}

\begin{tabular}{lll}
\multicolumn{3}{l}{\textbf{MIT PRÄPOSITIONEN}}\\
& \textbf{sich schminken für} + Akk & (to put on makeup for)\\
& \textbf{sich schminken mit} + Dat & (to put makeup on with)\\
\end{tabular}
\hfill
\begin{tabular}{lll}
\multicolumn{3}{l}{\textbf{TRENNBARE VERBEN}}\\
& \textbf{ab$|$schminken} & (to remove makeup)\\
& \textbf{an$|$schminken} & (to apply makeup)\\
\end{tabular}
\vspace{-5pt}
\end{verbentry}

% sich schützen (to protect oneself)
\begin{verbentry}{
  style = {default},
  toc = {sich schützen (to protect oneself)},
  name = {sich schützen},
  english = {to protect oneself},
  type = {regular},
  auxiliary = {haben}
}
 
    
     \vspace{5pt}

    \hrule
    
    \vspace{5pt}

  \begin{tabular}{rl}
     \multicolumn{2}{l}{\textbf{PRÄSENS} }\\
    ich & \textbf{schütze mich}\\
    du & \textbf{schützt dich}\\
    er/sie/es & \textbf{schützt sich}\\
    wir & \textbf{schützen uns}\\
    ihr & \textbf{schützt euch}\\
    sie/Sie & \textbf{schützen sich}\\
  \end{tabular}
  \hfill
     \begin{tabular}{rl}
    \multicolumn{2}{l}{\textbf{PRÄTERITUM} }\\
    ich & \textbf{schützte mich}\\
    du & \textbf{schütztest dich}\\
    er/sie/es & \textbf{schützte sich}\\
    wir & \textbf{schützten uns}\\
    ihr & \textbf{schütztet euch}\\
    sie/Sie & \textbf{schützten sich}\\
  \end{tabular}
  \hfill
  \begin{tabular}{rl}
     \multicolumn{2}{l}{\textbf{IMPERATIV} }\\
   & \textbf{schütz(e) dich!}\\
   & \textbf{schützen wir uns!}\\
   & \textbf{schützt euch!}\\
   & \textbf{schützen Sie sich!}\\
   & \vspace{10pt}
  \end{tabular}

    \vspace{10pt}

 \begin{tabular}{rl}
     \multicolumn{2}{l}{\textbf{PERFEKT}}\\
   & haben + sich + \textbf{geschützt}\\
  \end{tabular}
\hfill
   \begin{tabular}{rl}
    \multicolumn{2}{l}{\textbf{FUTURE I} }\\
   & werden + sich + \textbf{schützen}\\
  \end{tabular}
\hfill
   \begin{tabular}{rl}
      \multicolumn{2}{l}{\textbf{PLUSQUAMPERFEKT}}\\
   & hatten + sich + \textbf{geschützt}\\
  \end{tabular}

   \vspace{10pt}

   \begin{tabular}{lll}
      \multicolumn{3}{l}{\textbf{MIT PRÄPOSITIONEN}}\\
& \textbf{sich schützen vor} + Dat & (to protect oneself from)\\
& \textbf{sich schützen gegen} + Akk & (to protect oneself against)\\
\end{tabular}
  \hfill
   \begin{tabular}{lll}
\multicolumn{3}{l}{\textbf{TRENNBARE VERBEN}}\\
& \textbf{---} & \\
\end{tabular}

  \vspace{-5pt}
\end{verbentry}

% sich sehnen (to long for)
\begin{verbentry}{
  style = {default},
  toc = {sich sehnen (to long for)},
  name = {sich sehnen},
  english = {to long for},
  type = {regular, reflexive},
  auxiliary = {haben}
}
 
    
     \vspace{5pt}

    \hrule
    
    \vspace{5pt}

  \begin{tabular}{rl}
     \multicolumn{2}{l}{\textbf{PRÄSENS} }\\
    ich & \textbf{sehne mich}\\
    du & \textbf{sehnst dich}\\
    er/sie/es & \textbf{sehnt sich}\\
    wir & \textbf{sehnen uns}\\
    ihr & \textbf{sehnt euch}\\
    sie/Sie & \textbf{sehnen sich}\\
  \end{tabular}
  \hfill
     \begin{tabular}{rl}
    \multicolumn{2}{l}{\textbf{PRÄTERITUM} }\\
    ich & \textbf{sehnte mich}\\
    du & \textbf{sehntest dich}\\
    er/sie/es & \textbf{sehnte sich}\\
    wir & \textbf{sehnten uns}\\
    ihr & \textbf{sehntet euch}\\
    sie/Sie & \textbf{sehnten sich}\\
  \end{tabular}
  \hfill
  \begin{tabular}{rl}
     \multicolumn{2}{l}{\textbf{IMPERATIV} }\\
   & \textbf{sehn(e) dich!}\\
   & \textbf{sehnen wir uns!}\\
   & \textbf{sehnt euch!}\\
   & \textbf{sehnen Sie sich!}\\
   & \vspace{10pt}
  \end{tabular}

    \vspace{10pt}

 \begin{tabular}{rl}
     \multicolumn{2}{l}{\textbf{PERFEKT}}\\
   & haben + sich + \textbf{gesehnt}\\
  \end{tabular}
\hfill
   \begin{tabular}{rl}
    \multicolumn{2}{l}{\textbf{FUTURE I} }\\
   & werden + sich + \textbf{sehnen}\\
  \end{tabular}
\hfill
   \begin{tabular}{rl}
      \multicolumn{2}{l}{\textbf{PLUSQUAMPERFEKT}}\\
   & hatten + sich + \textbf{gesehnt}\\
  \end{tabular}

   \vspace{10pt}

   \begin{tabular}{lll}
      \multicolumn{3}{l}{\textbf{MIT PRÄPOSITIONEN}}\\
& \textbf{sich sehnen nach} + Dat & (to long for)\\
\end{tabular}
  \hfill
   \begin{tabular}{lll}
\multicolumn{3}{l}{\textbf{TRENNBARE VERBEN}}\\
& \textbf{---} & \\
\end{tabular}

  \vspace{-5pt}
\end{verbentry}

% sich setzen (to sit down / take a seat)
\begin{verbentry}{
  style = {default},
  toc = {sich setzen (to sit down / take a seat)},
  name = {sich setzen},
  english = {to sit down / take a seat},
  type = {regular},
  auxiliary = {haben}
}


\vspace{5pt}
\hrule
\vspace{5pt}

\begin{tabular}{rl}
\multicolumn{2}{l}{\textbf{PRÄSENS}}\\
ich & \textbf{setze mich}\\
du & \textbf{setzt dich}\\
er/sie/es & \textbf{setzt sich}\\
wir & \textbf{setzen uns}\\
ihr & \textbf{setzt euch}\\
sie/Sie & \textbf{setzen sich}\\
\end{tabular}
\hfill
\begin{tabular}{rl}
\multicolumn{2}{l}{\textbf{PRÄTERITUM}}\\
ich & \textbf{setzte mich}\\
du & \textbf{setztest dich}\\
er/sie/es & \textbf{setzte sich}\\
wir & \textbf{setzten uns}\\
ihr & \textbf{setztet euch}\\
sie/Sie & \textbf{setzten sich}\\
\end{tabular}
\hfill
\begin{tabular}{rl}
\multicolumn{2}{l}{\textbf{IMPERATIV}}\\
& \textbf{setze dich!}\\
& \textbf{setzen wir uns!}\\
& \textbf{setzt euch!}\\
& \textbf{setzen Sie sich!}\\
\end{tabular}

\vspace{10pt}

\begin{tabular}{rl}
\multicolumn{2}{l}{\textbf{PERFEKT}}\\
& haben + sich + \textbf{gesetzt}\\
\end{tabular}
\hfill
\begin{tabular}{rl}
\multicolumn{2}{l}{\textbf{FUTURE I}}\\
& werden + sich + \textbf{setzen}\\
\end{tabular}
\hfill
\begin{tabular}{rl}
\multicolumn{2}{l}{\textbf{PLUSQUAMPERFEKT}}\\
& hatten + sich + \textbf{gesetzt}\\
\end{tabular}

\vspace{10pt}

\begin{tabular}{lll}
\multicolumn{3}{l}{\textbf{MIT PRÄPOSITIONEN}}\\
& \textbf{sich setzen auf} + Akk & (to sit down on)\\
& \textbf{sich setzen neben} + Akk & (to sit next to)\\
\end{tabular}
\hfill
\begin{tabular}{lll}
\multicolumn{3}{l}{\textbf{TRENNBARE VERBEN}}\\
& \textbf{hin$|$setzen} & (to sit down)\\
\end{tabular}
\vspace{-5pt}
\end{verbentry}

% sich sorgen (to worry)
\begin{verbentry}{
  style = {default},
  toc = {sich sorgen (to worry)},
  name = {sich sorgen},
  english = {to worry},
  type = {regular},
  auxiliary = {haben}
}
 
    
       \vspace{5pt}

    \hrule
    
    \vspace{5pt}

  \begin{tabular}{rl}
     \multicolumn{2}{l}{\textbf{PRÄSENS} }\\
    ich & \textbf{sorge mich}\\
    du & \textbf{sorgst dich}\\
    er/sie/es & \textbf{sorgt sich}\\
    wir & \textbf{sorgen uns}\\
    ihr & \textbf{sorgt euch}\\
    sie/Sie & \textbf{sorgen sich}\\
  \end{tabular}
  \hfill
     \begin{tabular}{rl}
    \multicolumn{2}{l}{\textbf{PRÄTERITUM} }\\
    ich & \textbf{sorgte mich}\\
    du & \textbf{sorgtest dich}\\
    er/sie/es & \textbf{sorgte sich}\\
    wir & \textbf{sorgten uns}\\
    ihr & \textbf{sorgtet euch}\\
    sie/Sie & \textbf{sorgten sich}\\
  \end{tabular}
  \hfill
  \begin{tabular}{rl}
     \multicolumn{2}{l}{\textbf{IMPERATIV} }\\
   & \textbf{sorg(e) (du) dich!}\\
   & \textbf{sorgen wir uns!}\\
   & \textbf{sorgt (ihr) euch!}\\
   & \textbf{sorgen Sie sich!}\\
   & \vspace{10pt}
  \end{tabular}

    \vspace{10pt}

 \begin{tabular}{rl}
     \multicolumn{2}{l}{\textbf{PERFEKT}}\\
    &haben + sich + \textbf{gesorgt}\\
  \end{tabular}
\hfill
   \begin{tabular}{rl}
   
    \multicolumn{2}{l}{\textbf{FUTURE I} }\\
    &werden + sich + \textbf{sorgen}\\
 
  \end{tabular}
\hfill
   \begin{tabular}{rl}
      \multicolumn{2}{l}{\textbf{PLUSQUAMPERFEKT}}\\
    &hatten + sich + \textbf{gesorgt}\\
  \end{tabular}
  
\vspace{10pt}

   \begin{tabular}{lll}
      \multicolumn{3}{l}{\textbf{MIT PRÄPOSITIONEN}}\\
& \textbf{sich sorgen um} + Akk & (to worry about something)
  \end{tabular}
  \hfill
   \begin{tabular}{lll}
\multicolumn{3}{l}{\textbf{TRENNBARE VERBEN}}\\
& \textbf{---} & \\
\end{tabular}
  
  \vspace{-5pt}
\end{verbentry}

% sich stellen (to surrender)
\begin{verbentry}{
  style = {acc},
  toc = {sich stellen (to surrender)},
  name = {sich stellen},
  english = {to surrender},
  type = {regular, + Akkusativ},
  auxiliary = {haben}
}
 
    
       \vspace{5pt}

    \hrule
    
    \vspace{5pt}

  \begin{tabular}{rl}
     \multicolumn{2}{l}{\textbf{PRÄSENS} }\\
    ich & \textbf{stelle mich}\\
    du & \textbf{stellst dich}\\
    er/sie/es & \textbf{stellt sich}\\
    wir & \textbf{stellen uns}\\
    ihr & \textbf{stellt euch}\\
    sie/Sie & \textbf{stellen sich}\\
  \end{tabular}
  \hfill
     \begin{tabular}{rl}
    \multicolumn{2}{l}{\textbf{PRÄTERITUM} }\\
    ich & \textbf{stellte mich}\\
    du & \textbf{stelltest dich}\\
    er/sie/es & \textbf{stellte sich}\\
    wir & \textbf{stellten uns}\\
    ihr & \textbf{stelltet euch}\\
    sie/Sie & \textbf{stellten sich}\\
  \end{tabular}
  \hfill
  \begin{tabular}{rl}
     \multicolumn{2}{l}{\textbf{IMPERATIV} }\\
   & \textbf{stell(e) (du) dich!}\\
   & \textbf{stellen wir uns!}\\
   & \textbf{stellt (ihr) euch!}\\
   & \textbf{stellen Sie sich!}\\
   & \vspace{10pt}
  \end{tabular}

    \vspace{10pt}

 \begin{tabular}{rl}
     \multicolumn{2}{l}{\textbf{PERFEKT}}\\
    &haben + sich + \textbf{gestellt}\\
  \end{tabular}
\hfill
   \begin{tabular}{rl}
   
    \multicolumn{2}{l}{\textbf{FUTURE I} }\\
    &werden + sich + \textbf{stellen}\\
 
  \end{tabular}
\hfill
   \begin{tabular}{rl}
      \multicolumn{2}{l}{\textbf{PLUSQUAMPERFEKT}}\\
    &hatten + sich + \textbf{gestellt}\\
  \end{tabular}
  
\vspace{10pt}

   \begin{tabular}{lll}
      \multicolumn{3}{l}{\textbf{MIT PRÄPOSITIONEN}}\\
\vspace{0pt}
  \end{tabular}
  \hfill
   \begin{tabular}{lll}
      \multicolumn{3}{l}{\textbf{TRENNBARE VERBEN}}\\
& \textbf{sich vor$|$stellen} & (to introduce oneself)
  \end{tabular}
  
  \vspace{-5pt}
\end{verbentry}

% sich teilen (to share)
\begin{verbentry}{
  style = {dat},
  toc = {sich teilen (to share)},
  name = {sich teilen},
  english = {to share},
  type = {regular, + Dativ},
  auxiliary = {haben}
}
 
    
     \vspace{5pt}

    \hrule
    
    \vspace{5pt}

  \begin{tabular}{rl}
     \multicolumn{2}{l}{\textbf{PRÄSENS} }\\
    ich & \textbf{teile mir}\\
    du & \textbf{teilst dir}\\
    er/sie/es & \textbf{teilt sich}\\
    wir & \textbf{teilen uns}\\
    ihr & \textbf{teilt euch}\\
    sie/Sie & \textbf{teilen sich}\\
  \end{tabular}
  \hfill
     \begin{tabular}{rl}
    \multicolumn{2}{l}{\textbf{PRÄTERITUM} }\\
    ich & \textbf{teilte mir}\\
    du & \textbf{teiltest dir}\\
    er/sie/es & \textbf{teilte sich}\\
    wir & \textbf{teilten uns}\\
    ihr & \textbf{teiltet euch}\\
    sie/Sie & \textbf{teilten sich}\\
  \end{tabular}
  \hfill
  \begin{tabular}{rl}
     \multicolumn{2}{l}{\textbf{IMPERATIV} }\\
   & \textbf{teil(e) (du) dir!}\\
   & \textbf{teilen wir uns!}\\
   & \textbf{teilt (ihr) euch!}\\
   & \textbf{teilen Sie sich!}\\
   & \vspace{10pt}
  \end{tabular}

    \vspace{10pt}

 \begin{tabular}{rl}
     \multicolumn{2}{l}{\textbf{PERFEKT}}\\
    &haben + sich + \textbf{geteilt}\\
  \end{tabular}
\hfill
   \begin{tabular}{rl}
   
    \multicolumn{2}{l}{\textbf{FUTURE I} }\\
    &werden + sich + \textbf{teilen}\\
 
  \end{tabular}
\hfill
   \begin{tabular}{rl}
      \multicolumn{2}{l}{\textbf{PLUSQUAMPERFEKT}}\\
    &hatten + sich + \textbf{geteilt}\\
  \end{tabular}
  
\vspace{10pt}

   \begin{tabular}{lll}
      \multicolumn{3}{l}{\textbf{MIT PRÄPOSITIONEN}}\\
& \textbf{sich teilen in} + Akk & (to split into parts)\\
& \textbf{sich teilen mit} + Dat & (to share with someone)
  \end{tabular}
\hfill
\begin{tabular}{lll}
\multicolumn{3}{l}{\textbf{TRENNBARE VERBEN}}\\
& \textbf{---} & \\
\end{tabular}
 % \hfill
 %  \begin{tabular}{lll}
  %    \multicolumn{3}{l}{\textbf{TRENNBARE VERBEN}}\\
%\vspace{10pt}
 % \end{tabular}


  \vspace{-5pt}
\end{verbentry}

% sich treffen (to meet up)
\begin{verbentry}{
  style = {acc},
  toc = {sich treffen (to meet up)},
  name = {sich treffen},
  english = {to come forward},
  type = {irregular, + Akkusativ},
  auxiliary = {haben}
}
 
    
       \vspace{5pt}

    \hrule
    
    \vspace{5pt}

  \begin{tabular}{rl}
     \multicolumn{2}{l}{\textbf{PRÄSENS} }\\
    ich & \textbf{treffe mich}\\
    du & \textbf{triffst dich}\\
    er/sie/es & \textbf{trifft sich}\\
    wir & \textbf{treffen uns}\\
    ihr & \textbf{trefft euch}\\
    sie/Sie & \textbf{treffen sich}\\
  \end{tabular}
  \hfill
     \begin{tabular}{rl}
    \multicolumn{2}{l}{\textbf{PRÄTERITUM} }\\
    ich & \textbf{traf mich}\\
    du & \textbf{trafst dich}\\
    er/sie/es & \textbf{traf sich}\\
    wir & \textbf{trafen uns}\\
    ihr & \textbf{traft euch}\\
    sie/Sie & \textbf{trafen sich}\\
  \end{tabular}
  \hfill
  \begin{tabular}{rl}
     \multicolumn{2}{l}{\textbf{IMPERATIV} }\\
   & \textbf{triff(e) (du) dich!}\\
   & \textbf{treffen wir uns!}\\
   & \textbf{trefft (ihr) euch!}\\
   & \textbf{treffen Sie sich!}\\
   & \vspace{10pt}
  \end{tabular}

    \vspace{10pt}

 \begin{tabular}{rl}
     \multicolumn{2}{l}{\textbf{PERFEKT}}\\
    &haben + sich + \textbf{getroffen}\\
  \end{tabular}
\hfill
   \begin{tabular}{rl}
   
    \multicolumn{2}{l}{\textbf{FUTURE I} }\\
    &werden + sich + \textbf{treffen}\\
 
  \end{tabular}
\hfill
   \begin{tabular}{rl}
      \multicolumn{2}{l}{\textbf{PLUSQUAMPERFEKT}}\\
    &hatten + sich + \textbf{getroffen}\\
  \end{tabular}
  
\vspace{10pt}

   \begin{tabular}{lll}
      \multicolumn{3}{l}{\textbf{MIT PRÄPOSITIONEN}}\\
& \textbf{sich treffen mit} + Dat & (to meet someone)\\
& \textbf{sich treffen zu} + Dat & (to meet for a purpose)\\
  \end{tabular}
  \hfill
   \begin{tabular}{lll}
\multicolumn{3}{l}{\textbf{TRENNBARE VERBEN}}\\
& \textbf{---} & \\
\end{tabular}


 
  
  \vspace{-5pt}
\end{verbentry}

% sich trennen (to separate)
\begin{verbentry}{
  style = {acc},
  toc = {sich trennen (to separate)},
  name = {sich trennen},
  english = {to separate},
  type = {regular, + Akkusativ},
  auxiliary = {haben}
}
 
    
       \vspace{5pt}

    \hrule
    
    \vspace{5pt}

  \begin{tabular}{rl}
     \multicolumn{2}{l}{\textbf{PRÄSENS} }\\
    ich & \textbf{trenne mich}\\
    du & \textbf{trennst dich}\\
    er/sie/es & \textbf{trennt sich}\\
    wir & \textbf{trennen uns}\\
    ihr & \textbf{trennt euch}\\
    sie/Sie & \textbf{trennen sich}\\
  \end{tabular}
  \hfill
     \begin{tabular}{rl}
    \multicolumn{2}{l}{\textbf{PRÄTERITUM} }\\
    ich & \textbf{trennte mich}\\
    du & \textbf{trenntest dich}\\
    er/sie/es & \textbf{trennte sich}\\
    wir & \textbf{trennten uns}\\
    ihr & \textbf{trenntet euch}\\
    sie/Sie & \textbf{trennten sich}\\
  \end{tabular}
  \hfill
  \begin{tabular}{rl}
     \multicolumn{2}{l}{\textbf{IMPERATIV} }\\
   & \textbf{trenn(e) (du) dich!}\\
   & \textbf{trennen wir uns!}\\
   & \textbf{trennt (ihr) euch!}\\
   & \textbf{trennen Sie sich!}\\
   & \vspace{10pt}
  \end{tabular}

    \vspace{10pt}

 \begin{tabular}{rl}
     \multicolumn{2}{l}{\textbf{PERFEKT}}\\
    &haben + sich + \textbf{getrennt}\\
  \end{tabular}
\hfill
   \begin{tabular}{rl}
   
    \multicolumn{2}{l}{\textbf{FUTURE I} }\\
    &werden + sich + \textbf{trennen}\\
 
  \end{tabular}
\hfill
   \begin{tabular}{rl}
      \multicolumn{2}{l}{\textbf{PLUSQUAMPERFEKT}}\\
    &hatten + sich + \textbf{getrennt}\\
  \end{tabular}
  
\vspace{10pt}

   \begin{tabular}{lll}
      \multicolumn{3}{l}{\textbf{MIT PRÄPOSITIONEN}}\\
& \textbf{sich trennen von} + Dat & (to separate from someone / something)
  \end{tabular}
\hfill
\begin{tabular}{lll}
\multicolumn{3}{l}{\textbf{TRENNBARE VERBEN}}\\
& \textbf{---} & \\
\end{tabular}
 % \hfill
 %  \begin{tabular}{lll}
  %    \multicolumn{3}{l}{\textbf{TRENNBARE VERBEN}}\\
%\vspace{10pt}
 % \end{tabular}
  
  \vspace{-5pt}
\end{verbentry}

% sich unterhalten (to chat)
\begin{verbentry}{
  style = {default},
  toc = {sich unterhalten (to chat)},
  name = {sich unterhalten},
  english = {to chat},
  type = {irregular},
  auxiliary = {haben}
}
 
    
       \vspace{5pt}

    \hrule
    
    \vspace{5pt}

  \begin{tabular}{rl}
     \multicolumn{2}{l}{\textbf{PRÄSENS} }\\
    ich & \textbf{unterhalte mich}\\
    du & \textbf{unterhältst dich}\\
    er/sie/es & \textbf{unterhält sich}\\
    wir & \textbf{unterhalten uns}\\
    ihr & \textbf{unterhaltet euch}\\
    sie/Sie & \textbf{unterhalten sich}\\
  \end{tabular}
  \hfill
     \begin{tabular}{rl}
    \multicolumn{2}{l}{\textbf{PRÄTERITUM} }\\
    ich & \textbf{underhielt mich}\\
    du & \textbf{underhieltst dich}\\
    er/sie/es & \textbf{underhielt sich}\\
    wir & \textbf{underhielten uns}\\
    ihr & \textbf{underhieltet euch}\\
    sie/Sie & \textbf{underhielten sich}\\
  \end{tabular}
  \hfill
  \begin{tabular}{rl}
     \multicolumn{2}{l}{\textbf{IMPERATIV} }\\
   & \textbf{unterhalt(e) (du) dich!}\\
   & \textbf{unterhalten wir uns!}\\
   & \textbf{unterhaltt (ihr) euch!}\\
   & \textbf{unterhalten Sie sich!}\\
   & \vspace{10pt}
  \end{tabular}

    \vspace{10pt}

 \begin{tabular}{rl}
     \multicolumn{2}{l}{\textbf{PERFEKT}}\\
    &haben + sich + \textbf{unterhalten}\\
  \end{tabular}
\hfill
   \begin{tabular}{rl}
   
    \multicolumn{2}{l}{\textbf{FUTURE I} }\\
    &werden + sich + \textbf{unterhalten}\\
 
  \end{tabular}
\hfill
   \begin{tabular}{rl}
      \multicolumn{2}{l}{\textbf{PLUSQUAMPERFEKT}}\\
    &hatten + sich + \textbf{unterhalten}\\
  \end{tabular}
  
\vspace{10pt}

   \begin{tabular}{lll}
      \multicolumn{3}{l}{\textbf{MIT PRÄPOSITIONEN}}\\
& \textbf{sich unterhalten über} + Akk & (to chat about)\\
& \textbf{sich unterhalten mit} + Dat & (to chat mit)\\
  \end{tabular}
 \hfill
   \begin{tabular}{lll}
\multicolumn{3}{l}{\textbf{TRENNBARE VERBEN}}\\
& \textbf{---} & \\
\end{tabular}
  
  \vspace{-5pt}
\end{verbentry}

% sich verabreden (to arrange / make an appointment)
\begin{verbentry}{
  style = {default},
  toc = {sich verabreden (to arrange / make an appointment)},
  name = {sich verabreden},
  english = {to arrange / make an appointment},
  type = {regular},
  auxiliary = {haben}
}


\vspace{5pt}
\hrule
\vspace{5pt}

\begin{tabular}{rl}
\multicolumn{2}{l}{\textbf{PRÄSENS}}\\
ich & \textbf{verabrede mich}\\
du & \textbf{verabredest dich}\\
er/sie/es & \textbf{verabredet sich}\\
wir & \textbf{verabreden uns}\\
ihr & \textbf{verabredet euch}\\
sie/Sie & \textbf{verabreden sich}\\
\end{tabular}
\hfill
\begin{tabular}{rl}
\multicolumn{2}{l}{\textbf{PRÄTERITUM}}\\
ich & \textbf{verabredete mich}\\
du & \textbf{verabredetest dich}\\
er/sie/es & \textbf{verabredete sich}\\
wir & \textbf{verabredeten uns}\\
ihr & \textbf{verabredetet euch}\\
sie/Sie & \textbf{verabredeten sich}\\
\end{tabular}
\hfill
\begin{tabular}{rl}
\multicolumn{2}{l}{\textbf{IMPERATIV}}\\
& \textbf{verabrede dich!}\\
& \textbf{verabreden wir uns!}\\
& \textbf{verabredet euch!}\\
& \textbf{verabreden Sie sich!}\\
\end{tabular}

\vspace{10pt}

\begin{tabular}{rl}
\multicolumn{2}{l}{\textbf{PERFEKT}}\\
& haben + sich + \textbf{verabredet}\\
\end{tabular}
\hfill
\begin{tabular}{rl}
\multicolumn{2}{l}{\textbf{FUTURE I}}\\
& werden + sich + \textbf{verabreden}\\
\end{tabular}
\hfill
\begin{tabular}{rl}
\multicolumn{2}{l}{\textbf{PLUSQUAMPERFEKT}}\\
& hatten + sich + \textbf{verabredet}\\
\end{tabular}

\vspace{10pt}

\begin{tabular}{lll}
\multicolumn{3}{l}{\textbf{MIT PRÄPOSITIONEN}}\\
& \textbf{sich verabreden mit} + Dat & (to arrange with)\\
\end{tabular}
\hfill
\begin{tabular}{lll}
\multicolumn{3}{l}{\textbf{TRENNBARE VERBEN}}\\
& \textbf{---} & \\
\end{tabular}
\vspace{-5pt}
\end{verbentry}

% sich verlassen (to rely)
\begin{verbentry}{
  style = {default},
  toc = {sich verlassen (to rely)},
  name = {sich verlassen},
  english = {to rely},
  type = {irregular},
  auxiliary = {haben}
}
 
    
     \vspace{5pt}

    \hrule
    
    \vspace{5pt}

  \begin{tabular}{rl}
     \multicolumn{2}{l}{\textbf{PRÄSENS} }\\
    ich & \textbf{verlasse mich}\\
    du & \textbf{verlässt dich}\\
    er/sie/es & \textbf{verlässt sich}\\
    wir & \textbf{verlassen uns}\\
    ihr & \textbf{verlasst euch}\\
    sie/Sie & \textbf{verlassen sich}\\
  \end{tabular}
  \hfill
     \begin{tabular}{rl}
    \multicolumn{2}{l}{\textbf{PRÄTERITUM} }\\
    ich & \textbf{verließ mich}\\
    du & \textbf{verließest dich}\\
    er/sie/es & \textbf{verließ sich}\\
    wir & \textbf{verließen uns}\\
    ihr & \textbf{verließt euch}\\
    sie/Sie & \textbf{verließen sich}\\
  \end{tabular}
  \hfill
  \begin{tabular}{rl}
     \multicolumn{2}{l}{\textbf{IMPERATIV} }\\
   & \textbf{verlass(e) dich!}\\
   & \textbf{verlassen wir uns!}\\
   & \textbf{verlasst euch!}\\
   & \textbf{verlassen Sie sich!}\\
   & \vspace{10pt}
  \end{tabular}

    \vspace{10pt}

 \begin{tabular}{rl}
     \multicolumn{2}{l}{\textbf{PERFEKT}}\\
   & haben + sich + \textbf{verlassen}\\
  \end{tabular}
\hfill
   \begin{tabular}{rl}
    \multicolumn{2}{l}{\textbf{FUTURE I} }\\
   & werden + sich + \textbf{verlassen}\\
  \end{tabular}
\hfill
   \begin{tabular}{rl}
      \multicolumn{2}{l}{\textbf{PLUSQUAMPERFEKT}}\\
   & hatten + sich + \textbf{verlassen}\\
  \end{tabular}

   \vspace{10pt}

   \begin{tabular}{lll}
      \multicolumn{3}{l}{\textbf{MIT PRÄPOSITIONEN}}\\
& \textbf{sich verlassen auf} + Akk & (to rely on)\\
\end{tabular}
  \hfill
   \begin{tabular}{lll}
\multicolumn{3}{l}{\textbf{TRENNBARE VERBEN}}\\
& \textbf{---} & \\
\end{tabular}

  \vspace{-5pt}
\end{verbentry}

% sich verletzen (to injure oneself)
\begin{verbentry}{
  style = {default},
  toc = {sich verletzen (to injure oneself)},
  name = {sich verletzen},
  english = {to injure oneself},
  type = {regular},
  auxiliary = {haben}
}


\vspace{5pt}
\hrule
\vspace{5pt}

\begin{tabular}{rl}
\multicolumn{2}{l}{\textbf{PRÄSENS}}\\
ich & \textbf{verletze mich}\\
du & \textbf{verletzt dich}\\
er/sie/es & \textbf{verletzt sich}\\
wir & \textbf{verletzen uns}\\
ihr & \textbf{verletzt euch}\\
sie/Sie & \textbf{verletzen sich}\\
\end{tabular}
\hfill
\begin{tabular}{rl}
\multicolumn{2}{l}{\textbf{PRÄTERITUM}}\\
ich & \textbf{verletzte mich}\\
du & \textbf{verletztest dich}\\
er/sie/es & \textbf{verletzte sich}\\
wir & \textbf{verletzten uns}\\
ihr & \textbf{verletz tet euch}\\
sie/Sie & \textbf{verletzten sich}\\
\end{tabular}
\hfill
\begin{tabular}{rl}
\multicolumn{2}{l}{\textbf{IMPERATIV}}\\
& \textbf{verletze dich (du) nicht!}\\
& \textbf{verletzen wir uns nicht!}\\
& \textbf{verletzt euch (ihr) nicht!}\\
& \textbf{verletzen Sie sich nicht!}\\
\end{tabular}

\vspace{10pt}

\begin{tabular}{rl}
\multicolumn{2}{l}{\textbf{PERFEKT}}\\
& haben + sich + \textbf{verletzt}\\
\end{tabular}
\hfill
\begin{tabular}{rl}
\multicolumn{2}{l}{\textbf{FUTURE I}}\\
& werden + sich + \textbf{verletzen}\\
\end{tabular}
\hfill
\begin{tabular}{rl}
\multicolumn{2}{l}{\textbf{PLUSQUAMPERFEKT}}\\
& hatten + sich + \textbf{verletzt}\\
\end{tabular}

\vspace{10pt}

\begin{tabular}{lll}
\multicolumn{3}{l}{\textbf{MIT PRÄPOSITIONEN}}\\
& \textbf{sich verletzen an} + Dat & (to injure oneself on)\\
& \textbf{sich verletzen bei} + Dat & (to injure oneself while doing)\\
\end{tabular}
\hfill
\begin{tabular}{lll}
\multicolumn{3}{l}{\textbf{TRENNBARE VERBEN}}\\
& \textbf{---} & \\
\end{tabular}
\vspace{-5pt}
\end{verbentry}

% sich verlieben (to fall in love)
\begin{verbentry}{
  style = {default},
  toc = {sich verlieben (to fall in love)},
  name = {sich verlieben},
  english = {to fall in love},
  type = {regular},
  auxiliary = {haben}
}
 
    
       \vspace{5pt}

    \hrule
    
    \vspace{5pt}

  \begin{tabular}{rl}
     \multicolumn{2}{l}{\textbf{PRÄSENS} }\\
    ich & \textbf{verliebe mir}\\
    du & \textbf{verliebst dir}\\
    er/sie/es & \textbf{verliebt sich}\\
    wir & \textbf{verlieben uns}\\
    ihr & \textbf{verliebt euch}\\
    sie/Sie & \textbf{verlieben sich}\\
  \end{tabular}
  \hfill
     \begin{tabular}{rl}
    \multicolumn{2}{l}{\textbf{PRÄTERITUM} }\\
    ich & \textbf{verliebte mir}\\
    du & \textbf{verliebtest dir}\\
    er/sie/es & \textbf{verliebte sich}\\
    wir & \textbf{verliebten uns}\\
    ihr & \textbf{verliebtet euch}\\
    sie/Sie & \textbf{verliebten sich}\\
  \end{tabular}
  \hfill
  \begin{tabular}{rl}
     \multicolumn{2}{l}{\textbf{IMPERATIV} }\\
   & \textbf{verlieb(e) (du) dir!}\\
   & \textbf{verlieben wir uns!}\\
   & \textbf{verliebt (ihr) euch!}\\
   & \textbf{verlieben Sie sich!}\\
   & \vspace{10pt}
  \end{tabular}

    \vspace{10pt}

 \begin{tabular}{rl}
     \multicolumn{2}{l}{\textbf{PERFEKT}}\\
    &haben + sich + \textbf{verliebt}\\
  \end{tabular}
\hfill
   \begin{tabular}{rl}
   
    \multicolumn{2}{l}{\textbf{FUTURE I} }\\
    &werden + sich + \textbf{verlieben}\\
 
  \end{tabular}
\hfill
   \begin{tabular}{rl}
      \multicolumn{2}{l}{\textbf{PLUSQUAMPERFEKT}}\\
    &hatten + sich + \textbf{verliebt}\\
  \end{tabular}
  
\vspace{10pt}

   \begin{tabular}{lll}
      \multicolumn{3}{l}{\textbf{MIT PRÄPOSITIONEN}}\\
& \textbf{sich verlieben in} + Akk & (to fall in love with someone)
  \end{tabular}
  \hfill
   \begin{tabular}{lll}
\multicolumn{3}{l}{\textbf{TRENNBARE VERBEN}}\\
& \textbf{---} & \\
\end{tabular}
  
  \vspace{-5pt}
\end{verbentry}

% sich verschreiben (to commit / write incorrectly)
\begin{verbentry}{
  style = {default},
  toc = {sich verschreiben (to commit / write incorrectly)},
  name = {sich verschreiben},
  english = {to commit / write incorrectly},
  type = {irregular},
  auxiliary = {haben}
}


\vspace{5pt}
\hrule
\vspace{5pt}

\begin{tabular}{rl}
\multicolumn{2}{l}{\textbf{PRÄSENS}}\\
ich & \textbf{verschreibe mich}\\
du & \textbf{verschreibst dich}\\
er/sie/es & \textbf{verschreibt sich}\\
wir & \textbf{verschreiben uns}\\
ihr & \textbf{verschreibt euch}\\
sie/Sie & \textbf{verschreiben sich}\\
\end{tabular}
\hfill
\begin{tabular}{rl}
\multicolumn{2}{l}{\textbf{PRÄTERITUM}}\\
ich & \textbf{verschrieb mich}\\
du & \textbf{verschriebst dich}\\
er/sie/es & \textbf{verschrieb sich}\\
wir & \textbf{verschrieben uns}\\
ihr & \textbf{verschriebt euch}\\
sie/Sie & \textbf{verschrieben sich}\\
\end{tabular}
\hfill
\begin{tabular}{rl}
\multicolumn{2}{l}{\textbf{IMPERATIV}}\\
& \textbf{verschreibe (du) dich!}\\
& \textbf{verschreiben wir uns!}\\
& \textbf{verschreibt (ihr) euch!}\\
& \textbf{verschreiben Sie sich!}\\
\end{tabular}

\vspace{10pt}

\begin{tabular}{rl}
\multicolumn{2}{l}{\textbf{PERFEKT}}\\
&haben + sich + \textbf{verschrieben}\\
\end{tabular}
\hfill
\begin{tabular}{rl}
\multicolumn{2}{l}{\textbf{FUTURE I}}\\
& werden + sich + \textbf{verschreiben}\\
\end{tabular}
\hfill
\begin{tabular}{rl}
\multicolumn{2}{l}{\textbf{PLUSQUAMPERFEKT}}\\
& hatten + sich + \textbf{verschrieben}\\
\end{tabular}

\vspace{10pt}

\begin{tabular}{lll}
\multicolumn{3}{l}{\textbf{MIT PRÄPOSITIONEN}}\\
& \textbf{---} & \\
\end{tabular}
\hfill
\begin{tabular}{lll}
\multicolumn{3}{l}{\textbf{TRENNBARE VERBEN}}\\
& \textbf{---} & \\
\end{tabular}
\vspace{-5pt}
\end{verbentry}

% sich verstehen (to get along)
\begin{verbentry}{
  style = {acc},
  toc = {sich verstehen (to get along)},
  name = {sich verstehen},
  english = {to get along},
  type = {irregular, + Akkusativ},
  auxiliary = {haben}
}
 
    
       \vspace{5pt}

    \hrule
    
    \vspace{5pt}

  \begin{tabular}{rl}
     \multicolumn{2}{l}{\textbf{PRÄSENS} }\\
    ich & \textbf{verstehe mich}\\
    du & \textbf{verstehst dich}\\
    er/sie/es & \textbf{versteht sich}\\
    wir & \textbf{verstehen uns}\\
    ihr & \textbf{versteht euch}\\
    sie/Sie & \textbf{verstehen sich}\\
  \end{tabular}
  \hfill
     \begin{tabular}{rl}
    \multicolumn{2}{l}{\textbf{PRÄTERITUM} }\\
    ich & \textbf{verstand mich}\\
    du & \textbf{verstandest dich}\\
    er/sie/es & \textbf{verstand sich}\\
    wir & \textbf{verstanden uns}\\
    ihr & \textbf{verstandet euch}\\
    sie/Sie & \textbf{verstanden sich}\\
  \end{tabular}
  \hfill
  \begin{tabular}{rl}
     \multicolumn{2}{l}{\textbf{IMPERATIV} }\\
   & \textbf{versteh(e) (du) dich!}\\
   & \textbf{verstehen wir uns!}\\
   & \textbf{versteht (ihr) euch!}\\
   & \textbf{verstehen Sie sich!}\\
   & \vspace{10pt}
  \end{tabular}

    \vspace{10pt}

 \begin{tabular}{rl}
     \multicolumn{2}{l}{\textbf{PERFEKT}}\\
    &haben + sich + \textbf{verstanden}\\
  \end{tabular}
\hfill
   \begin{tabular}{rl}
   
    \multicolumn{2}{l}{\textbf{FUTURE I} }\\
    &werden + sich + \textbf{verstehen}\\
 
  \end{tabular}
\hfill
   \begin{tabular}{rl}
      \multicolumn{2}{l}{\textbf{PLUSQUAMPERFEKT}}\\
    &hatten + sich + \textbf{verstanden}\\
  \end{tabular}
  
\vspace{10pt}

   \begin{tabular}{lll}
      \multicolumn{3}{l}{\textbf{MIT PRÄPOSITIONEN}}\\
& \textbf{sich verlieben mit} + Dat & (to get along with)
  \end{tabular}
\hfill
\begin{tabular}{lll}
\multicolumn{3}{l}{\textbf{TRENNBARE VERBEN}}\\
& \textbf{---} & \\
\end{tabular}
 % \hfill
 %  \begin{tabular}{lll}
  %    \multicolumn{3}{l}{\textbf{TRENNBARE VERBEN}}\\
%\vspace{10pt}
 % \end{tabular}
  
  \vspace{-5pt}
\end{verbentry}

% sich verwählen (to misdial / miscall)
\begin{verbentry}{
  style = {default},
  toc = {sich verwählen (to misdial / miscall)},
  name = {sich verwählen},
  english = {to misdial / miscall},
  type = {regular},
  auxiliary = {haben}
}


\vspace{5pt}
\hrule
\vspace{5pt}

\begin{tabular}{rl}
\multicolumn{2}{l}{\textbf{PRÄSENS}}\\
ich & \textbf{verwähle mich}\\
du & \textbf{verwählst dich}\\
er/sie/es & \textbf{verwählt sich}\\
wir & \textbf{verwählen uns}\\
ihr & \textbf{verwählt euch}\\
sie/Sie & \textbf{verwählen sich}\\
\end{tabular}
\hfill
\begin{tabular}{rl}
\multicolumn{2}{l}{\textbf{PRÄTERITUM}}\\
ich & \textbf{verwählte mich}\\
du & \textbf{verwähltest dich}\\
er/sie/es & \textbf{verwählte sich}\\
wir & \textbf{verwählten uns}\\
ihr & \textbf{verwähltet euch}\\
sie/Sie & \textbf{verwählten sich}\\
\end{tabular}
\hfill
\begin{tabular}{rl}
\multicolumn{2}{l}{\textbf{IMPERATIV}}\\
& \textbf{verwähle dich!}\\
& \textbf{verwählen wir uns!}\\
& \textbf{verwählt euch!}\\
& \textbf{verwählen Sie sich!}\\
\end{tabular}

\vspace{10pt}

\begin{tabular}{rl}
\multicolumn{2}{l}{\textbf{PERFEKT}}\\
& haben + sich + \textbf{verwählt}\\
\end{tabular}
\hfill
\begin{tabular}{rl}
\multicolumn{2}{l}{\textbf{FUTURE I}}\\
& werden + sich + \textbf{verwählen}\\
\end{tabular}
\hfill
\begin{tabular}{rl}
\multicolumn{2}{l}{\textbf{PLUSQUAMPERFEKT}}\\
& hatten + sich + \textbf{verwählt}\\
\end{tabular}

\vspace{10pt}

\begin{tabular}{lll}
\multicolumn{3}{l}{\textbf{MIT PRÄPOSITIONEN}}\\
& \textbf{---} & \\
\end{tabular}
\hfill
\begin{tabular}{lll}
\multicolumn{3}{l}{\textbf{TRENNBARE VERBEN}}\\
& \textbf{---} & \\
\end{tabular}
\vspace{-5pt}
\end{verbentry}

% sich waschen (to wash oneself)
\begin{verbentry}{
  style = {acc},
  toc = {sich waschen (to wash oneself)},
  name = {sich waschen},
  english = {to wash},
  type = {irregular, + Akkusativ},
  auxiliary = {haben}
}
 
    
       \vspace{5pt}

    \hrule
    
    \vspace{5pt}

  \begin{tabular}{rl}
     \multicolumn{2}{l}{\textbf{PRÄSENS} }\\
    ich & \textbf{wasche mich}\\
    du & \textbf{wäschst dich}\\
    er/sie/es & \textbf{wäscht sich}\\
    wir & \textbf{waschen uns}\\
    ihr & \textbf{wascht euch}\\
    sie/Sie & \textbf{waschen sich}\\
  \end{tabular}
  \hfill
     \begin{tabular}{rl}
    \multicolumn{2}{l}{\textbf{PRÄTERITUM} }\\
    ich & \textbf{wunch mich}\\
    du & \textbf{wunchst dich}\\
    er/sie/es & \textbf{wunch sich}\\
    wir & \textbf{wunchen uns}\\
    ihr & \textbf{wuncht euch}\\
    sie/Sie & \textbf{wunchen sich}\\
  \end{tabular}
  \hfill
  \begin{tabular}{rl}
     \multicolumn{2}{l}{\textbf{IMPERATIV} }\\
   & \textbf{wasch(e) (du) dich!}\\
   & \textbf{waschen wir uns!}\\
   & \textbf{wascht (ihr) euch!}\\
   & \textbf{waschen Sie sich!}\\
   & \vspace{10pt}
  \end{tabular}

    \vspace{10pt}

 \begin{tabular}{rl}
     \multicolumn{2}{l}{\textbf{PERFEKT}}\\
    &haben + sich + \textbf{gewaschen}\\
  \end{tabular}
\hfill
   \begin{tabular}{rl}
   
    \multicolumn{2}{l}{\textbf{FUTURE I} }\\
    &werden + sich + \textbf{waschen}\\
 
  \end{tabular}
\hfill
   \begin{tabular}{rl}
      \multicolumn{2}{l}{\textbf{PLUSQUAMPERFEKT}}\\
    &hatten + sich + \textbf{gewaschen}\\
  \end{tabular}
  
\vspace{10pt}

\begin{tabular}{lll}
\multicolumn{3}{l}{\textbf{MIT PRÄPOSITIONEN}}\\
& \textbf{---} & \\
\end{tabular}
\hfill
\begin{tabular}{lll}
\multicolumn{3}{l}{\textbf{TRENNBARE VERBEN}}\\
& \textbf{---} & \\
\end{tabular}

\vspace{-5pt}
\end{verbentry}

% sich wundern (to wonder / be surprised)
\begin{verbentry}{
  style = {default},
  toc = {sich wundern (to wonder / be surprised)},
  name = {sich wundern},
  english = {to wonder / be suprised},
  type = {regular},
  auxiliary = {haben}
}
 
    
       \vspace{5pt}

    \hrule
    
    \vspace{5pt}

  \begin{tabular}{rl}
     \multicolumn{2}{l}{\textbf{PRÄSENS} }\\
    ich & \textbf{wundere mich}\\
    du & \textbf{wunderst dich}\\
    er/sie/es & \textbf{wundert sich}\\
    wir & \textbf{wundern uns}\\
    ihr & \textbf{wundert euch}\\
    sie/Sie & \textbf{wundern sich}\\
  \end{tabular}
  \hfill
     \begin{tabular}{rl}
    \multicolumn{2}{l}{\textbf{PRÄTERITUM} }\\
    ich & \textbf{wunderte mich}\\
    du & \textbf{wundertest dich}\\
    er/sie/es & \textbf{wunderte sich}\\
    wir & \textbf{wunderten uns}\\
    ihr & \textbf{wundertet euch}\\
    sie/Sie & \textbf{wunderten sich}\\
  \end{tabular}
  \hfill
  \begin{tabular}{rl}
     \multicolumn{2}{l}{\textbf{IMPERATIV} }\\
   & \textbf{wunder(e) (du) dich!}\\
   & \textbf{wundern wir uns!}\\
   & \textbf{wundert (ihr) euch!}\\
   & \textbf{wundern Sie sich!}\\
   & \vspace{10pt}
  \end{tabular}

    \vspace{10pt}

 \begin{tabular}{rl}
     \multicolumn{2}{l}{\textbf{PERFEKT}}\\
    &haben + sich + \textbf{gewundert}\\
  \end{tabular}
\hfill
   \begin{tabular}{rl}
   
    \multicolumn{2}{l}{\textbf{FUTURE I} }\\
    &werden + sich + \textbf{wundern}\\
 
  \end{tabular}
\hfill
   \begin{tabular}{rl}
      \multicolumn{2}{l}{\textbf{PLUSQUAMPERFEKT}}\\
    &hatten + sich + \textbf{gewundert}\\
  \end{tabular}
  
\vspace{10pt}

   \begin{tabular}{lll}
      \multicolumn{3}{l}{\textbf{MIT PRÄPOSITIONEN}}\\
& \textbf{sich wundern über} + Akk & (to wonder / be suprised about)
  \end{tabular}
  \hfill
  \begin{tabular}{lll}
\multicolumn{3}{l}{\textbf{TRENNBARE VERBEN}}\\
& \textbf{---} & \\
\end{tabular}
  
  \vspace{-5pt}
\end{verbentry}

% sich ziehen (to drag)
\begin{verbentry}{
  style = {acc},
  toc = {sich ziehen (to drag)},
  name = {sich ziehen},
  english = {to drag},
  type = {irregular, + Akkusativ},
  auxiliary = {haben}
}
 
    
       \vspace{5pt}

    \hrule
    
    \vspace{5pt}

  \begin{tabular}{rl}
     \multicolumn{2}{l}{\textbf{PRÄSENS} }\\
    ich & \textbf{ziehe mich}\\
    du & \textbf{ziehst dich}\\
    er/sie/es & \textbf{zieht sich}\\
    wir & \textbf{ziehen uns}\\
    ihr & \textbf{zieht euch}\\
    sie/Sie & \textbf{ziehen sich}\\
  \end{tabular}
  \hfill
     \begin{tabular}{rl}
    \multicolumn{2}{l}{\textbf{PRÄTERITUM} }\\
    ich & \textbf{zog mich}\\
    du & \textbf{zogst dich}\\
    er/sie/es & \textbf{zog sich}\\
    wir & \textbf{zogen uns}\\
    ihr & \textbf{zogt euch}\\
    sie/Sie & \textbf{zogen sich}\\
  \end{tabular}
  \hfill
  \begin{tabular}{rl}
     \multicolumn{2}{l}{\textbf{IMPERATIV} }\\
   & \textbf{zieh(e) (du) dich!}\\
   & \textbf{ziehen wir uns!}\\
   & \textbf{zieht (ihr) euch!}\\
   & \textbf{ziehen Sie sich!}\\
   & \vspace{10pt}
  \end{tabular}

    \vspace{10pt}

 \begin{tabular}{rl}
     \multicolumn{2}{l}{\textbf{PERFEKT}}\\
    &haben + sich + \textbf{gezogen}\\
  \end{tabular}
\hfill
   \begin{tabular}{rl}
   
    \multicolumn{2}{l}{\textbf{FUTURE I} }\\
    &werden + sich + \textbf{ziehen}\\
 
  \end{tabular}
\hfill
   \begin{tabular}{rl}
      \multicolumn{2}{l}{\textbf{PLUSQUAMPERFEKT}}\\
    &hatten + sich + \textbf{gezogen}\\
  \end{tabular}
  
\vspace{10pt}

   \begin{tabular}{lll}
      \multicolumn{3}{l}{\textbf{MIT PRÄPOSITIONEN}}\\
\vspace{0pt}
  \end{tabular}
  \hfill
   \begin{tabular}{lll}
    \multicolumn{3}{l}{\textbf{TRENNBARE VERBEN}}\\
& \textbf{sich an$|$ziehan} + Akk & (to get dressed)
 \end{tabular}
  
  \vspace{-5pt}
\end{verbentry}
